% STATUS: reviewed
\chapter{API}\label{ch:api}

% One route of a table: the methods and the path on a line of their own, across both columns;
% the row under it says who may call and what the route does.
\providecommand{\apiroute}[2]{\multicolumn{2}{@{}l@{}}{\route{#1}{#2}}\\*}

Šis skyrius yra sistemos HTTP sąsajos žinynas: kiekvienas maršrutas, kas gali jį kviesti, ką jis
priima ir ką atmeta. Skyrių sudaro penkios dalys. \Cref{sec:api:conventions} pateikia visiems
maršrutams bendras taisykles --- kur keliauja užklausa, prieigos raktas, klaidos apvalkalas,
puslapiavimas, ribos. \Cref{sec:api:routes} išvardija API (Django) maršrutus pagal sritis, kiekvienai
sričiai po lentelę. \Cref{sec:api:assistant} pateikia tris viešus asistento paslaugos maršrutus,
kuriais vyksta pokalbis. \Cref{sec:api:internal} aprašo vidines duris, kurias dvi paslaugos atveria
viena kitai ir niekam kitam. \Cref{sec:api:client} surenka, ką API kviečianti programa turi žinoti
apie darbą, trunkantį ilgiau nei užklausa, ir apie failus. \term{Klientas} šiame skyriuje --- bet
kuri programa, kviečianti API: konsolė, naudotojo vardu veikianti asistento paslauga arba scenarijus.

Skyrius sako, ką daro kiekvienas maršrutas, bet ne kodėl: priežastys išdėstytos skyriuje apie
maršruto dalyką, nurodytame prie kiekvienos lentelės. API failas \file{urls.py} yra sąsajos turinys
--- kiekvienas API maršrutas jame užima vieną eilutę, --- o kiekvienas rodinių modulis prasideda
savo maršrutų ir jų užklausų turinio sąrašu.


% ============================================================
\section{Bendros taisyklės}\label{sec:api:conventions}
% ============================================================

\subsection{Trys kelių šeimos}\label{sec:api:families}

Sistema turi vieną viešą adresą. Už jo stovi \term{įeinančio srauto tarpinis serveris} (toliau ---
įėjimas), Caddy serveris \code{legalops-endpoint}, sukonfigūruotas faile \file{endpoint/Caddyfile}.
Jis klauso 80~prievado --- vienintelio, kurį diegimas paskelbia, --- ir kiekvieną užklausą
paskirsto pagal jos kelią; laimi pirmas atitikmuo (\cref{fig:api:routing}):

\begin{enumerate}
  \item \code{/api/assistant/*} keliauja į asistento paslaugą (\code{legalops-assistant},
    3000 prievadas). Tai trys pokalbio maršrutai (\cref{sec:api:assistant}).
  \item \code{/api/*} keliauja į API (\code{legalops-django}, 8000 prievadas): tai visi kiti šio
    skyriaus maršrutai, kuriuos klientas gali pasiekti.
  \item Visa kita keliauja į konsolės ir viešosios svetainės žiniatinklio serverį
    (\code{legalops-vite}), kuris atsako failais, niekada ne duomenimis (\cref{ch:console}).
\end{enumerate}

Prieš šiuos tris dar stovi vienas kelias: \code{/dbgate/*} veda į DBGate, duomenų bazės
administravimo įrankį. Prieš perduodamas tokią užklausą, įėjimas paklausia API, ar kviečiantysis
yra administratorius (\route{GET}{/api/auth/checkauth/admin}), o atsisakymą grąžina tokį, koks
jis yra (\cref{ch:ops}).

\begin{figure}[H]
\centering
\begin{tikzpicture}[node distance=4mm and 12mm]
  \node[svc, minimum width=44mm] (assistant) {asistento paslauga\\{\scriptsize\code{legalops-assistant:3000}}};
  \node[svc, minimum width=44mm, below=of assistant] (django) {API (Django)\\{\scriptsize\code{legalops-django:8000}}};
  \node[svc, minimum width=44mm, below=of django] (vite) {konsolė ir viešoji svetainė\\{\scriptsize\code{legalops-vite:80}}};
  \node[svc, minimum width=44mm, below=of vite] (dbgate) {DBGate\\{\scriptsize\code{legalops-dbgate:3000}}};
  \node[svc, minimum height=16mm, anchor=east] (in) at ([xshift=-36mm]$(django.west)!.5!(vite.west)$) {įėjimas\\{\scriptsize\code{legalops-endpoint}}\\{\scriptsize 80 prievadas}};
  \node[ext, left=8mm of in] (client) {naršyklė\\ar scenarijus};
  \node[note, below=3mm of in, text width=34mm] {prieš \code{/dbgate/*} užklausą jis paklausia API, ar kviečiantysis yra administratorius};
  \draw[flow] (client) -- (in);
  \draw[flow] (in.east) -- ++(4mm,0) |- (assistant.west) node[lbl, pos=.8] {\code{/api/assistant/*}};
  \draw[flow] (in.east) -- ++(4mm,0) |- (django.west) node[lbl, pos=.8] {\code{/api/*}};
  \draw[flow] (in.east) -- ++(4mm,0) |- (vite.west) node[lbl, pos=.8] {visa kita};
  \draw[flow] (in.east) -- ++(4mm,0) |- (dbgate.west) node[lbl, pos=.8] {\code{/dbgate/*}};
  \begin{scope}[on background layer]
    \node[lane, fit=(assistant)(django)(vite)(dbgate)] (lane) {};
  \end{scope}
  \node[lanelabel, anchor=south west] at (lane.north west) {izoliuotas tinklas: jokio savo prievado};
\end{tikzpicture}
\caption{Kur keliauja užklausa. Įėjimas yra vienintelė paslauga su paskelbtu prievadu; paslaugą jis
renkasi vien pagal kelią. Niekas, kas yra po \code{/internal/} ar \code{/health}, nepersiunčiama nei
į API, nei į asistento paslaugą.}\label{fig:api:routing}
\end{figure}

Trys įėjimo veiksmai liečia kiekvieną klientą. Užklausai, kurios antraštė \code{Origin} nurodo kitą
svetainę, atsakoma 403 su tuščiu turiniu, dar jai nepasiekus jokios paslaugos. Įėjimas perduoda
gautą antraštę \code{X-Forwarded-For} ir nukopijuoja ją į \code{X-Real-IP}; pagal šią antraštę API
skaičiuoja prisijungimo bandymus ir rašo veiksmų žurnalo eilutes (\code{client_ip} faile
\file{common/http.py}). Reikšmė yra tokia, kokią atnešė gaunamoji užklausa; jos nesant API
naudoja ryšio adresą, t.~y. paties įėjimo adresą. O kai už kelio stovinti paslauga neveikia arba
paleidžiama iš naujo, įėjimas pats atsako 502 su tuščiu turiniu. Saugumo antraštės, kurias jis
prideda prie kiekvieno atsakymo, išskyrus savo paties 403, aprašytos \cref{ch:security}.

\subsection{Užklausos ir atsakymai}\label{sec:api:json}

API kalba tik JSON: užklausos turinys yra JSON objektas, o kiekvienas atsakymas --- ir klaidos
atsakymas --- yra JSON, kurį rašo \code{json_response} faile \file{common/http.py}, UTF-8 simbolius
paliekant tokius, kokie yra. Išimčių nedaug, ir jos įvardytos lentelėse: du įkėlimai priima
\code{multipart/form-data}, trys eksportai atsako failu, vienas maršrutas atsako SVG brėžiniu, o
pokalbis atsako srautu.

\begin{description}[style=unboxed, leftmargin=0pt]
  \item[Keliai.] Joks kelias nesibaigia pasviruoju brūkšniu (\code{APPEND_SLASH} išjungtas); tas
    pats kelias su juo duoda 404. Nežinomas kelias atsako 404 \code{NOT_FOUND}, o metodas, kurio
    maršrutas neaptarnauja, --- 405 \code{METHOD_NOT_ALLOWED} su antrašte \code{Allow}; abu kaip
    JSON, niekada ne kaip HTML puslapis.
  \item[Laukų pavadinimai.] camelCase (\code{displayName}, \code{companyId}). Skiriasi dvi vietos:
    pokalbio istorijos eilutės išlaiko jas skaitančios pokalbio vykdymo aplinkos pavadinimus
    (\code{parent_id}), o vidinės durys naudoja snake\_case (\code{user_id}, \code{input_tokens}).
  \item[Reikšmės.] Data ir laikas --- ISO~8601 su poslinkiu, pvz., \code{2026-09-27T18:00:00+00:00};
    data --- \code{YYYY-MM-DD}; kaina --- JAV dolerių skaičius.
  \item[Identifikatoriai.] Kiekviena maršruto sukurta eilutė gauna uuid4 simbolių eilutę kaip
    identifikatorių. Į instrukcijas kreipiamasi pagal jų identifikatorių (\en{slug}), į
    instrukcijos versiją --- pagal jos numerį. Identifikatorius, atkeliavęs užklausos turinyje ar
    parametruose, priimamas tik kaip netuščia ne ilgesnė nei 64 simbolių eilutė (\code{str_id});
    visa kita laikoma identifikatoriaus nebuvimu, todėl objektas ar sąrašas niekada nepasiekia
    paieškos duomenų bazėje.
  \item[Transakcijos.] Kiekviena užklausa vykdoma vienoje duomenų bazės transakcijoje
    (\code{ATOMIC_REQUESTS} faile \file{settings.py}): nepavykusi užklausa nepalieka nieko
    įrašyto iki pusės. Už jos ribų lieka keturi maršrutai, laukiantys kažko lėto, --- būklės
    patikra, dokumento įkėlimas, bandomasis instrukcijos paleidimas ir MCP ryšio patikra.
\end{description}

\subsection{Prieigos raktas ir kas gali kviesti}\label{sec:api:auth}

Klientas prisijungia per \route{POST}{/api/auth/login} ir gautą prieigos raktą siunčia kiekvienoje
vėlesnėje užklausoje kaip \code{Authorization: Bearer <token>}. Prieigos raktas yra neskaidri
atsitiktinė eilutė (256 atsitiktiniai bitai; duomenų bazė saugo tik jos SHA-256), galiojanti
\code{SESSION_DAYS}~=~30 dienų; paskyra išlaiko \code{SESSIONS_PER_USER}~=~10 naujausių sesijų,
todėl vienuoliktasis prisijungimas užbaigia seniausią. Paskyra įkeliama iš naujo per kiekvieną
užklausą, todėl išjungta paskyra užrakinama iš karto. Savarankiškos registracijos ir autentifikacijos
slapukais nėra: vienintelį API nustatomą slapuką \code{compliance_gate} skaito tik vartai prieš
DBGate. Užklausai be galiojančio prieigos rakto atsakoma 401 \code{AUTH_REQUIRED}. Mechanizmas yra
faile \file{users/auth.py}; kodėl jis taip sukurtas, aiškinama \cref{ch:security}.

Yra dvi rolės: \code{user} ir \code{admin}. \Cref{sec:api:routes} lentelės nurodo, kas gali kviesti
maršrutą, vienu iš keturių žodžių:

\begin{table}[H]
\caption{Kas gali kviesti maršrutą --- maršrutų lentelėse vartojami žodžiai}\label{tab:api:who}
\begin{tabularx}{\textwidth}{@{}L{22mm}Y@{}}
\toprule
\hd{Žodis} & \hd{Reikšmė} \\
\midrule
viešas & Prieigos rakto nereikia. Vieši tik būklės patikra ir prisijungimas. \\
prisijungęs & Bet kuri aktyvi paskyra su galiojančiu prieigos raktu (\code{@require_auth}). \\
savininkas & Bet kuris prisijungęs naudotojas --- savo eilutėms; administratorius --- visoms
  eilutėms. Kurios eilutės yra naudotojo, nustato \term{nuosavybės taisyklė}, aprašyta toliau. \\
administratorius & Tik administratoriai (\code{@require_role("admin")} arba rolės patikra rodinio viduje).
  Naudotojas gauna 403 \code{FORBIDDEN}. \\
\bottomrule
\end{tabularx}
\end{table}

Nuosavybės taisyklė užrašyta vieną kartą kiekvienai eilučių rūšiai --- kaip funkcija, grąžinanti
eilutes, kurias kviečiantysis gali liesti, --- ir kiekvienas tos rūšies maršrutas pradeda nuo jos:

\begin{itemize}
  \item \emph{Įmonė} priklauso ją sukūrusiam naudotojui (\code{visible_companies} faile
    \file{companies/api/company_views.py}). Jos leidimai, taikomi standartai, rolės, ataskaitos,
    išvados, veiksmai ir programos dokumentai seka paskui įmonę.
  \item \emph{Dokumentas} priklauso jį įkėlusiam naudotojui (\code{visible_documents}), nesvarbu,
    ar jis susietas su įmone.
  \item \emph{Užduotis} priklauso ją sukėlusiam naudotojui (\code{_visible} faile
    \file{jobs/api/views.py}), o \emph{modelio iškvietimas} DI sąnaudų žurnale --- naudotojui,
    kuriam jis priskaičiuotas.
  \item \emph{Pokalbis} priklauso savo naudotojui, ir čia administratorius nėra išimtis
    (\code{visible_threads} faile \file{chat/api/views.py}): per pokalbių maršrutus niekas
    neskaito kito žmogaus pokalbių.
  \item \emph{Teisės aktų bazę} ir pareigų registrą gali skaityti kiekvienas prisijungęs.
    Naudotojas mato tik įjungtus teisės aktus; administratorius mato visus. Pareigų maršrutai niekada
    neišvardija įmonės vidaus reikalavimų.
\end{itemize}

\paragraph{Ką reiškia 404.} Eilutė, kurios kviečiantysis negali matyti, jam neegzistuoja: maršrutas
atsako 404 \code{NOT_FOUND}, tą patį atsakymą kaip niekada neišduotam identifikatoriui, todėl API
niekada nepatvirtina, kad kieno nors kito įmonė, dokumentas, ataskaita ar pokalbis egzistuoja. 403
paliktas dviem atvejams, kai juo nieko neatskleidžiama: kai naudotojas kviečia tik
administratoriams skirtą maršrutą ir kai naudotojas DI sąnaudų maršrutų klausia kitų paskyrų
skaičių.

\subsection{Klaidos}\label{sec:api:errors}

Kiekviena API klaida turi vieną pavidalą --- \term{klaidos apvalkalą} (\code{json_error} faile
\file{common/http.py}):

\begin{codeblock}
{"error": "An analysis of this company is already running", "code": "ALREADY_RUNNING"}
\end{codeblock}

\code{error} yra angliškas sakinys žmogui; \code{code} --- pastovus pavadinimas, pagal kurį klientas
renkasi, ką daryti. HTTP būsenos kodas nurodo nesėkmės klasę. \Cref{tab:api:codes} išvardija visus
API grąžinamus kodus; maršrutų lentelės nurodo konkrečiam maršrutui būdingus kodus ir praleidžia
tuos, kuriuos gali grąžinti bet kuris maršrutas (\code{INVALID_BODY}, \code{AUTH_REQUIRED},
\code{FORBIDDEN}, \code{NOT_FOUND}, \code{METHOD_NOT_ALLOWED}).

\begin{table}[H]
\caption{API (Django) klaidų kodai pagal HTTP būseną}\label{tab:api:codes}
\begin{tabularx}{\textwidth}{@{}L{12mm}Y@{}}
\toprule
\hd{Būsena} & \hd{Kodai} \\
\midrule
400 & \code{INVALID_BODY} (turinyje nėra JSON objekto), \code{MISSING_FIELDS},
  \code{INVALID_FIELD}, \code{INVALID_QUERY} (užklausos parametras), \code{EMPTY_UPDATE} (nėra ką
  keisti), \code{MISSING_FILE}, \code{BAD_REQUEST} (NUL simbolis kelyje ar parametruose),
  \code{INVALID_VALUE} (reikšmė, kurios duomenų bazė negali priimti: už intervalo ribų arba netinkamai sudaryta), \code{INVALID_EMAIL},
  \code{INVALID_ROLE}, \code{WEAK_PASSWORD}, \code{WRONG_PASSWORD}, \code{LAST_ADMIN},
  \code{SELF_LOCKOUT}, \code{UNKNOWN_MODEL}, \code{NOTHING_TO_ASSESS}, \code{NOTE_REQUIRED} \\
401 & \code{AUTH_REQUIRED} (sesijos nėra arba ji nebegalioja), \code{INVALID_CREDENTIALS}
  (prisijungimas) \\
403 & \code{FORBIDDEN} (rolė), \code{ROLE_NOT_ALLOWED} (pokalbio žinutė, kuri nėra paties
  naudotojo), \code{INTERNAL_ONLY} (vidinės durys be paslapties) \\
404 & \code{NOT_FOUND}; vidinėse duryse dar \code{PROMPT_NOT_CONFIGURED} \\
405 & \code{METHOD_NOT_ALLOWED}, su antrašte \code{Allow} \\
409 & Užklausa sudaryta teisingai, bet eilutė nėra tam tinkamos būsenos:
  \code{ALREADY_RUNNING}, \code{STILL_RUNNING}, \code{NOT_CANCELLABLE}, \code{NOT_FINISHED},
  \code{NOT_WAITING}, \code{NOT_DRAFT}, \code{NOT_APPROVED}, \code{NOT_READY},
  \code{NEWER_VERSION}, \code{RULES_BROKEN}, \code{CLOSED_BY_REPORT}, \code{CORPUS_EMPTY},
  \code{NOT_INGESTED}, \code{NO_FILE}, \code{NO_DOCUMENTS}, \code{UNCHANGED},
  \code{UNKNOWN_KIND}; arba pavadinimas užimtas: \code{EMAIL_TAKEN}, \code{DUPLICATE_NAME},
  \code{SLUG_TAKEN}, \code{ALREADY_ADOPTED}, \code{MESSAGE_ID_TAKEN} \\
413 & \code{FILE_TOO_LARGE}, \code{QUOTA_EXCEEDED}, \code{MESSAGE_TOO_LARGE} \\
429 & \code{RATE_LIMITED} (su antrašte \code{Retry-After}), \code{BUDGET_EXCEEDED}
  (kviečiančiojo dienos biudžetas modelių iškvietimams) \\
500 & \code{SERVER_ERROR}; klaidos pėdsakas rašomas į žurnalą, niekada nesiunčiamas klientui.
  Reikšmė, kuri praėjo pačios API patikrą ir pažeidžia duomenų bazės apribojimą, --- išskyrus
  užimtą pavadinimą, kuriam atsakoma 409, --- taip pat baigiasi čia, ir sąmoningai: tai parodo
  tos patikros spragą, kurią 400 paslėptų \\
502 & \code{MODEL_ERROR} (bandomasis instrukcijos paleidimas, kurio modelio iškvietimas nepavyko) \\
503 & \code{MODELS_UNAVAILABLE} (nepavyksta perskaityti DI tinklų sąsajos modelių sąrašo) \\
\bottomrule
\end{tabularx}
\end{table}

Asistento paslauga savo klaidas apgaubia kitaip, nes tokio pavidalo tikisi konsolėje veikianti
pokalbio vykdymo aplinka (\code{errorMiddleware} asistento paslaugos faile
\file{middleware/errors.js}):

\begin{codeblock}
{"message": "Too many questions in a row — wait a moment",
 "error": {"code": "RATE_LIMITED", "message": "Too many questions in a row — wait a moment", "retryAfterS": 41}}
\end{codeblock}

\code{retryAfterS} pasirodo tik su 429 ir pakartojamas antraštėje \code{Retry-After}. Asistento
paslaugos kodai išvardyti \cref{sec:api:assistant}.

\subsection{Sąrašai: puslapiavimas ir filtrai}\label{sec:api:paging}

Sąrašo maršrutas atsako objektu, kuriame eilutės pateikiamos daugiskaitiniu pavadinimu
(\code{companies}, \code{reports}), o kai sąrašas puslapiuojamas, --- ir visų atitinkančių eilučių
skaičius \code{total}. Puslapiuojami tik ilgieji sąrašai; kiti grąžina viską iki nustatytos ribos
(\cref{tab:api:paging}). Naudojami du skaičių skaitytuvai, ir netinkamą reikšmę jie traktuoja
skirtingai: \code{query_int} reikšmę už jos intervalo ribų atmeta su 400 \code{INVALID_QUERY}, o
užduočių, veiksmų žurnalo ir nuostatų maršrutai ją priderina prie intervalo ir atmeta tik tai, kas
nėra skaičius.

\begin{table}[H]
\caption{Kaip puslapiuojami sąrašų maršrutai}\label{tab:api:paging}
\begin{tabularx}{\textwidth}{@{}L{55mm}L{27mm}Y@{}}
\toprule
\hd{Sąrašas} & \hd{\code{limit}: numatytoji, didžiausia} & \hd{Puslapiavimas} \\
\midrule
\code{/api/documents} & 500, 500 & \code{limit}, \code{offset}; grąžina \code{total}; už ribų: 400 \\
\code{/api/obligations} & 500, 2000 & \code{limit}, \code{offset}; grąžina \code{total}; už ribų: 400 \\
\code{/api/usage/calls} & 50, 200 & \code{limit}, \code{offset}; grąžina \code{total}; už ribų: 400 \\
\code{/api/corpus/documents/{id}/units} & 200, 1000 & \code{limit}, \code{offset}, priderinami; grąžina \code{total}, \code{offset}, \code{limit} \\
\code{/api/jobs} & 100, 500 & tik \code{limit}, priderinamas; naujausios pirmos \\
\code{/api/audit} & 100, 500 & žymeklis: \code{before} ir \code{beforeId}, paimti iš ankstesnio puslapio \code{next} \\
\code{/api/documents/{id}/text} & 20\,000, 100\,000 & simbolių langas: \code{offset}, \code{limit}; grąžina \code{total} \\
įmonės, pokalbiai, ataskaitos, programos dokumentai, veiksmai & --- & nepuslapiuojami; daugiausia 500, 300, 200, 300 ir 500 eilučių \\
\bottomrule
\end{tabularx}
\end{table}

Filtrai yra užklausos parametrai, pavadinti pagal lauką, kurį jie filtruoja (\code{companyId},
\code{status}, \code{kind}, \code{jurisdiction}). \code{q} visur, kur pasitaiko, reiškia teksto
paiešką. \code{scope=own} susiaurina administratoriaus įmonių ar dokumentų sąrašą iki paties
administratoriaus eilučių --- būtent to prašo pokalbių asistentas; \code{scope=all} ir
\code{userId} DI sąnaudų maršrutuose, priešingai, praplečia administratoriaus vaizdą. Filtro
reikšmę, kurios nėra jo sąraše, teisės aktų bazės ir pareigų sąrašai ignoruoja, o užduočių ir DI
sąnaudų sąrašai atmeta su 400 \code{INVALID_QUERY}; ataskaitų sąraše nežinoma \code{status}
reikšmė nieko neatitinka.

\subsection{Kalbos}\label{sec:api:languages}

Pati API turi vieną kalbą: jos klaidų sakiniai angliški, o klientas, rodantis savo formuluotes,
jas parenka pagal \code{code}. Kai maršrutas kuria tekstą skaitytojui, užklausa nurodo, kuria iš
dviejų produkto kalbų. Ataskaita ir programos dokumentas priima \code{language} --- \code{lt} arba
\code{en}, jo nesant lietuvių kalba --- ir išlaiko ją visam laikui. Veiksmų registro eksportas
\code{language} priima kaip užklausos parametrą. Pareigų išskyrimas priima
\code{language: "en"}, kitu atveju vyksta lietuviškai. Pokalbio ėjimas \code{locale} priima savo
turinyje; jo nesant sprendžia antraštė \code{Accept-Language}: anglų kalba, kai ji prasideda
\code{en}, ir lietuvių kalba visais kitais atvejais (\code{turnLanguage} asistento paslaugos faile
\file{routes/chat.js}). Teisės akto \code{language} teisės aktų bazėje yra faktas apie jo tekstą,
o ne išvesties pasirinkimas.

\subsection{Ribos}\label{sec:api:limits}

\paragraph{Užklausų dažnio ribos.} Dažniu ribojami du dalykai, abu su fiksuotu langu kiekvienam
raktui. Prisijungimas skaičiuoja bandymus kiekvienam kliento adresui --- \code{LOGIN_RATE_PER_IP},
numatytai 60, per \code{LOGIN_WINDOW_S}~=~60 sekundžių --- ir kiekvienai paskyrai:
\code{EMAIL_LIMIT}~=~10 bandymų per \code{EMAIL_WINDOW_S}~=~600 sekundžių; šiuo langu dalijasi ir
\route{POST}{/api/auth/change-password} slaptažodžio patikra. Langai yra lentelės
\code{rate_limit_windows} eilutės (\file{common/ratelimit.py}), todėl visi API procesai skaičiuoja
kartu, o atmestas bandymas vis tiek įskaitomas. Pokalbis leidžia \code{RATE_LIMIT_TURNS}~=~12
ėjimų vienam naudotojui per \code{RATE_LIMIT_WINDOW_MS}~=~60 sekundžių; jie skaičiuojami asistento
paslaugos atmintyje, ir ten atmestas ėjimas neįskaitomas. Abu atsako 429 \code{RATE_LIMITED} su
antrašte \code{Retry-After}. Joks kitas maršrutas dažniu neribojamas.

\paragraph{Dienos biudžetas.} Brangius maršrutus riboja pinigai, ne užklausų skaičius. Kiekvienas
naudotojas turi dienos biudžetą modelių iškvietimams JAV doleriais (\code{LLM_DAILY_BUDGET_USD},
numatytai 10,00~\$, arba paskyros \code{dailyBudgetUsd}; administratoriaus biudžetas neribotas,
jei nenustatytas). Pokalbio ėjimas ir bandomasis instrukcijos paleidimas atmetami su 429
\code{BUDGET_EXCEEDED}, kai nieko nebelieka; ataskaitos paleidimui, jos apimties patvirtinimui ir
programos dokumento paleidimui reikia, kad liktų \code{MIN_BUDGET_TO_START_USD}~=~0,25~\$.
Nuorodų tikrintuvas --- vienintelis maršrutas, kuris ne atsisako, o veikia supaprastintai.
\Cref{ch:models} aprašo, kaip iškvietimai įkainojami ir skaičiuojami.

\paragraph{Dydžiai.} Įėjimas neriboja jokio turinio; kiekviena paslauga taiko savo ribas. API
skaito turinį, kuris nėra įkeliamas failas, iki \code{DATA_UPLOAD_MAX_MEMORY_SIZE}~=~8~MiB
(\file{settings.py}); didesnį Django atmeta dar prieš tai, kai jį pamato bet kuris maršrutas.
Įmonės dokumentas gali būti iki 50~MiB (\code{MAX_FILE_BYTES}), o teisės aktų bazės failas ---
iki 60~MiB (\code{MAX_UPLOAD_BYTES}); viršijus įkėlimas atsako 413 \code{FILE_TOO_LARGE}.
Paskyra gali saugoti \code{UPLOAD_QUOTA_FILES}~=~1000 dokumentų ir iš viso
\code{UPLOAD_QUOTA_BYTES}~=~2~GiB; viršijus --- 413 \code{QUOTA_EXCEEDED}. Saugoma pokalbio žinutė
gali būti iki 400\,000 simbolių JSON pavidalu (413 \code{MESSAGE_TOO_LARGE}). Asistento paslauga
pokalbio ėjimui priima 512~kB JSON (413 \code{PAYLOAD_TOO_LARGE}). Teksto laukai turi savo ribas,
nurodytas lentelėse ten, kur jos svarbios.

\begin{codebox}{Griežtas JSON skaitymas}
Vienintelis užklausų turinio skaitytuvas, \code{get_json_object} faile \file{common/http.py},
priima JSON \emph{objektą} ir nieko daugiau. Jis atmeta, lyg turinys nebūtų išsiųstas:
\begin{itemize}
  \item masyvą ar paprastą reikšmę viršutiniame lygmenyje;
  \item \code{NaN} ir \code{Infinity}, užrašytus žodžiu ar gautus perpildžius (\code{1e999});
  \item NUL simbolį ar pavienį UTF-16 pakaitinį simbolį bet kuriame rakte ar eilutėje --- jų
    negali laikyti joks duomenų bazės tekstas;
  \item gilesnį nei \code{MAX_JSON_DEPTH}~=~32 įdėjimą.
\end{itemize}
Dauguma maršrutų tada atsako 400 \code{INVALID_BODY}; maršrutas, kurio turinys neprivalomas,
atmestą turinį skaito kaip tuščią. Laukai skaitomi tipizuotais skaitytuvais: \code{strict_bool}
priima JSON loginę reikšmę ir niekada eilutės \code{"false"}, \code{bounded_int} priima sveikąjį
skaičių savo intervale ir niekada eilutės. Maršrutui nežinomas raktas ignoruojamas. Prieš visus
maršrutus \code{RequestGuardMiddleware} (\file{common/middleware.py}) atsako 400
\code{BAD_REQUEST}, jei kelyje ar parametruose yra NUL simbolis, o reikšmę, kurios duomenų bazė
negali priimti, --- esančią už intervalo ribų ar netinkamai sudarytą, ---
paverčia 400 \code{INVALID_VALUE} vietoj 500. Kai reikšmė pažeidžia duomenų bazės apribojimą,
paliekamas 500 (\cref{tab:api:codes}).
\end{codebox}


% ============================================================
\section{API maršrutai}\label{sec:api:routes}
% ============================================================

\file{urls.py} turi 83 URL šablonus: 77 po \code{/api/}, išvardytus šiame poskyryje, ir 6 po
\code{/internal/assistant/}, išvardytus \cref{sec:api:internal}. Skaičiuojant pagal metodus, 83
šablonai aptarnauja 119 metodo ir kelio porų. Kiekvienoje toliau pateiktoje lentelėje vienas
įrašas skirtas vienam šablonui: pirmoje eilutėje metodai ir kelias, po jais --- kas gali kviesti
(\cref{tab:api:who}) ir ką maršrutas daro. Kelyje \code{{id}} žymi eilutės identifikatorių.
\enquote{Klaidos} nurodo tam maršrutui būdingus atsisakymus.

Pirmiausia du langeliai surūšiuoja maršrutus pagal klausimą, kurį šis dokumentas kelia
kiekvienai sistemos daliai: kas veikia.

\begin{aibox}{Maršrutai, kurie prieš atsakydami kviečia modelį}
Beveik kiekvienas maršrutas, kuriam reikia modelio, tik įtraukia užduotį į eilę ir atsako iš
karto; modelį paskui kviečia užduočių vykdytojas (\cref{ch:arch}). Penkios užklausų rūšys modelio
laukia pačios:
\begin{itemize}
  \item \route{POST}{/api/corpus/check-citations} --- vienas vertintojo iškvietimas nuorodų
    patikroje, kai užklausa nurodo, apie ką nuostata tariamai kalba (\cref{ch:company});
  \item \route{POST}{/api/prompts/{slug}/test} --- vienas iškvietimas su bandoma instrukcija
    (\cref{ch:prompts});
  \item \route{POST}{/api/corpus/search} ir \route{GET}{/api/documents/search} --- vienas
    užklausos vektorizavimas (\cref{ch:corpus});
  \item \route{POST}{/api/obligations} ir \code{PATCH}, keičiantis pareigos pavadinimą ar tekstą,
    --- vienas pareigos vektorizavimas (\cref{ch:obligations});
  \item \route{POST}{/api/assistant/chat} --- pats pokalbio ėjimas (\cref{ch:chat}).
\end{itemize}
\end{aibox}

\begin{humanbox}{Maršrutai, kuriais sprendžia žmogus}
Kiekvienas sprendimas, kurio sistema niekada nepriima pati, patenka per vieną maršrutą, pažymimas,
kas ir kada nusprendė, ir įrašomas į veiksmų žurnalą: ataskaitos apimties patvirtinimas
(\route{POST}{/api/reports/{id}/scope} su \code{approved: true}); išvados būsenos pakeitimas ir
jos priežasties įrašymas (\route{PATCH}{/api/reports/{id}/findings/{fid}}); ataskaitos
patvirtinimas (\route{PATCH}{/api/reports/{id}}); pareigos patvirtinimas, taisymas ar atmetimas
(\route{PATCH}{/api/obligations/{id}}, \route{POST}{/api/obligations/review}); programos dokumento
patvirtinimas (\route{PATCH}{/api/programme/{id}}); leidimo pažymėjimas patikrintu
(\route{PATCH}{/api/companies/{id}/credentials/{cid}} arba \code{POST}, kuris jį įrašo su
šia būsena); veiksmo uždarymas rankiniu būdu, su
pastaba (\route{PATCH}{/api/actions/{id}}).
\end{humanbox}

\subsection{Būklė}\label{sec:api:health}

Viena vieša patikra parodo, ar API gali dirbti savo darbą (\file{ops/api/views.py}).

\begin{xltabular}{\textwidth}{@{}L{19mm}Y@{}}
\caption{Maršrutai: būklė}\label{tab:api:health}\\
\toprule \hd{Kas} & \hd{Ką daro maršrutas} \\ \midrule \endfirsthead
\toprule \hd{Kas} & \hd{Ką daro maršrutas} \\ \midrule \endhead
\bottomrule \endlastfoot
\apiroute{GET HEAD}{/api/health}
viešas & Parengtis: duomenų bazė atsako į paprasčiausią užklausą, o į įkėlimų katalogą galima
  rašyti. 200 su \code{status: "ok"}, \code{service: "legalops-django"}, \code{database},
  \code{uploads}; arba 503 su \code{status: "error"} ir \code{reason} --- pirmąja nepavykusia
  patikra (\code{database} arba \code{uploads}); nurodoma, kuri patikra, bet niekada ne kodėl:
  priežastis yra žurnale. \\
\end{xltabular}

\subsection{Autentifikacija}\label{sec:api:authroutes}

Paties kviečiančiojo sesija (\file{users/api/auth_views.py}). \code{user} atsakyme yra paskyra
taip, kaip ją užrašo \code{serialize_user}: \code{id}, \code{email}, \code{displayName},
\code{role}, \code{active}, \code{dailyBudgetUsd}, \code{createdAt}, \code{lastLoginAt}.

\begin{xltabular}{\textwidth}{@{}L{19mm}Y@{}}
\caption{Maršrutai: autentifikacija}\label{tab:api:auth}\\
\toprule \hd{Kas} & \hd{Ką daro maršrutas} \\ \midrule \endfirsthead
\toprule \hd{Kas} & \hd{Ką daro maršrutas} \\ \midrule \endhead
\bottomrule \endlastfoot
\apiroute{POST}{/api/auth/login}
viešas & Turinys \code{email}, \code{password}. Grąžina \code{token} ir \code{user} ir nustato
  vartų slapuką. El.~pašto adresas lyginamas neatsižvelgiant į raidžių dydį. Nežinomas el.~pašto
  adresas, neteisingas slaptažodis ir išjungta paskyra gauna tą patį atsakymą, 401
  \code{INVALID_CREDENTIALS}, per tą patį laiką. Sėkmingi ir nesėkmingi bandymai įrašomi į veiksmų
  žurnalą. \emph{Klaidos:} 400 \code{MISSING_FIELDS}; 429 \code{RATE_LIMITED}. \\
\addlinespace
\apiroute{GET}{/api/auth/me}
prisijungęs & Paties kviečiančiojo \code{user}. Konsolė jį kviečia kiekvieną kartą pasileisdama;
  asistento paslauga jį kviečia, kad prieigos raktą paverstų naudotoju. \\
\addlinespace
\apiroute{GET}{/api/auth/checkauth/admin}
administratorius & Vartai, kurių įėjimas klausia prieš užklausą į \code{/dbgate/}. Sesija imama
  iš prieigos rakto arba --- kadangi naršyklėje atidarytas įrankis siunčia tik slapukus --- iš
  vartų slapuko. 200 su \code{user} aktyviam administratoriui; 401 \code{AUTH_REQUIRED}; 403
  \code{FORBIDDEN} bet kuriai kitai rolei. \\
\addlinespace
\apiroute{POST}{/api/auth/logout}
prisijungęs & Užbaigia pateikto prieigos rakto sesiją ir panaikina vartų slapuką. Kituose
  įrenginiuose naudotojas lieka prisijungęs. \\
\addlinespace
\apiroute{POST}{/api/auth/logout-all}
prisijungęs & Užbaigia visas kviečiančiojo sesijas, įskaitant pateiktąją. Įrašoma į veiksmų
  žurnalą. \\
\addlinespace
\apiroute{POST}{/api/auth/change-password}
prisijungęs & Turinys \code{currentPassword}, \code{newPassword}. Užbaigiamos visos \emph{kitos}
  paskyros sesijos. Slaptažodį sudaro mažiausiai 10 simbolių ir daugiausia 72 baitai.
  \emph{Klaidos:} 400 \code{WRONG_PASSWORD} (įrašoma į veiksmų žurnalą), \code{WEAK_PASSWORD}; 429
  \code{RATE_LIMITED}. \\
\end{xltabular}

\subsection{Naudotojai}\label{sec:api:users}

Paskyrų katalogas administratoriams (\file{users/api/users_views.py}). Kiekvienas pakeitimas
įrašomas į veiksmų žurnalą. Kiekvienam rolės ar būsenos pakeitimui ir šalinimui galioja dvi
apsaugos: paskutinio aktyvaus administratoriaus negalima nei pažeminti, nei išjungti, nei pašalinti
(400 \code{LAST_ADMIN}), ir niekas negali to padaryti savo paskyrai (400 \code{SELF_LOCKOUT}).

\begin{xltabular}{\textwidth}{@{}L{19mm}Y@{}}
\caption{Maršrutai: naudotojai}\label{tab:api:users}\\
\toprule \hd{Kas} & \hd{Ką daro maršrutas} \\ \midrule \endfirsthead
\toprule \hd{Kas} & \hd{Ką daro maršrutas} \\ \midrule \endhead
\bottomrule \endlastfoot
\apiroute{GET POST}{/api/users}
administratorius & \code{GET}: visos paskyros, įskaitant išjungtas, seniausios pirmos, kaip
  \code{users}. \code{POST}: turinys \code{email}, \code{displayName}, \code{password} ir
  \code{role} (jos nesant --- \code{user}); atsako 201 su paskyra. El.~pašto adresas saugomas
  mažosiomis raidėmis. \emph{Klaidos:} 400 \code{INVALID_EMAIL}, \code{MISSING_FIELDS},
  \code{INVALID_ROLE}, \code{WEAK_PASSWORD}; 409 \code{EMAIL_TAKEN} (nesvarbu, kokiomis raidėmis). \\
\addlinespace
\apiroute{GET PATCH DELETE}{/api/users/{id}}
administratorius & \code{PATCH}: bet kuris iš \code{displayName}, \code{role}, \code{active},
  \code{dailyBudgetUsd} (nuo 0 iki 10\,000; \code{null} reiškia numatytąjį, o 0 --- neribotą).
  \code{DELETE} išjungia paskyrą ir užbaigia jos sesijas; eilutės lieka. \code{DELETE} su
  \code{purge=true} visam laikui pašalina paskyrą kartu su jos įmonėmis, dokumentais ir
  pokalbiais ir grąžina, ką pašalino. \emph{Klaidos:} 400 \code{EMPTY_UPDATE},
  \code{INVALID_FIELD}; 409 \code{STILL_RUNNING} (šalinimas, kol kurios nors paskyros įmonės
  ataskaita yra eilėje ar vykdoma). \\
\addlinespace
\apiroute{POST}{/api/users/{id}/reset-password}
administratorius & Turinys \code{password}. Nustato paskyrai naują slaptažodį ir užbaigia visas
  jos sesijas. \emph{Klaidos:} 400 \code{WEAK_PASSWORD}. \\
\end{xltabular}

\subsection{Įmonės}\label{sec:api:companies}

Įmonės profilis ir tai, kas prie jo pririšta: leidimai, taikomi standartai ir rolės
(\file{companies/api/company_views.py}, \file{companies/api/standard_views.py},
\file{companies/api/role_views.py}). Ką šie laukai reiškia ataskaitai, aptariama \cref{ch:company}.
Kiekvienas keitimas įrašomas į veiksmų žurnalą.

\begin{xltabular}{\textwidth}{@{}L{19mm}Y@{}}
\caption{Maršrutai: įmonės, leidimai, standartai ir rolės}\label{tab:api:companies}\\
\toprule \hd{Kas} & \hd{Ką daro maršrutas} \\ \midrule \endfirsthead
\toprule \hd{Kas} & \hd{Ką daro maršrutas} \\ \midrule \endhead
\bottomrule \endlastfoot
\apiroute{GET POST}{/api/companies}
savininkas & \code{GET}: kviečiančiojo įmonės kaip \code{companies}, neseniausiai atnaujintos
  pirmos, daugiausia 500. \code{q} lyginamas žodis po žodžio su pavadinimu ir registracijos kodu;
  skyrybos ženklai, teisinės formos (UAB, LTD) ir žodis \enquote{company} praleidžiami.
  \code{scope=own} susiaurina administratoriaus sąrašą iki jo paties įmonių. \code{POST}:
  \code{name} ir bet kurie profilio laukai; atsako 201 su įmone; kviečiantysis tampa jos
  savininku. \emph{Klaidos:} 400 \code{MISSING_FIELDS}, \code{INVALID_FIELD}. \\
\addlinespace
\apiroute{GET PUT DELETE}{/api/companies/{id}}
savininkas & \code{GET}: profilis su \code{credentials}, \code{standards}, \code{roles} ir
  \code{documents}. \code{PUT} yra dalinis atnaujinimas: bet koks profilio laukų poaibis.
  \code{DELETE} pašalina įmonę su jos leidimais, taikomais standartais, rolėmis, ataskaitomis ir
  programos dokumentais; jos dokumentai atsiejami, o ne pašalinami. \emph{Klaidos:} 400
  \code{INVALID_FIELD}; 409 \code{STILL_RUNNING} (vykdoma ataskaita ar programos dokumentas). \\
\addlinespace
\apiroute{GET POST}{/api/companies/{id}/credentials}
savininkas & Leidimas --- tai įmonės turima licencija, sertifikatas, leidimas ar registracija.
  \code{POST}: \code{kind} (\code{certificate}, \code{license}, \code{permit},
  \code{registration}), \code{name} ir neprivalomai \code{issuer}, \code{identifier},
  \code{validFrom}, \code{validTo}, \code{scope}, \code{notes}, \code{status},
  \code{evidenceDocumentId}. Naujas leidimas yra \code{claimed}, nebent turinyje nurodyta
  \code{verified}. \emph{Klaidos:} 400 \code{MISSING_FIELDS}, \code{INVALID_FIELD}. \\
\addlinespace
\apiroute{PATCH DELETE}{/api/companies/{id}/credentials/{cid}}
savininkas & \code{PATCH}: bet koks tų pačių laukų poaibis. Nustačius \code{status} į
  \code{verified}, pažymima, kas ir kada patikrino; \code{claimed} žymą panaikina. Įrodymas turi
  būti vienas iš pačios įmonės dokumentų. \\
\addlinespace
\apiroute{GET POST}{/api/companies/{id}/standards}
savininkas & Įmonės taikomi standartai. \code{POST}: \code{documentId} (teisės aktų bazės
  standartas), \code{status} (\code{certified}, \code{implementing}, \code{planned}; jo nesant
  --- \code{implementing}), \code{scope}, \code{certificateId} (vienas iš įmonės leidimų),
  \code{notes}. \emph{Klaidos:} 400 \code{MISSING_FIELDS}, \code{INVALID_FIELD}; 409
  \code{ALREADY_ADOPTED}. \\
\addlinespace
\apiroute{PATCH DELETE}{/api/companies/{id}/standards/{sid}}
savininkas & \code{PATCH}: bet kuris iš \code{status}, \code{scope}, \code{certificateId}
  (\code{null} jį atsieja), \code{notes}. Pats standartas per šiuos maršrutus niekada nekeičiamas. \\
\addlinespace
\apiroute{GET POST}{/api/companies/{id}/roles}
savininkas & Rolė --- tai pareigybė įmonėje, pvz., atitikties pareigūnas, aprašyta pokalbių
  asistentui. \code{POST}: \code{name}, \code{context} (daugiausia 6\,000 simbolių),
  \code{sources} (daugiausia 50 įrašų iš \code{documentId}, \code{title}, \code{section}).
  Šaltinis, kuris identifikatoriumi arba pavadinimu nurodo vieną iš įmonės dokumentų,
  išsaugomas ir grąžinamas su abiem.
  \emph{Klaidos:} 400 \code{MISSING_FIELDS}, \code{INVALID_FIELD}; 409 \code{DUPLICATE_NAME}
  (toje pačioje įmonėje). \\
\addlinespace
\apiroute{GET PATCH DELETE}{/api/companies/{id}/roles/{rid}}
savininkas & \code{PATCH}: bet kuris iš \code{name}, \code{context}, \code{sources}. Pašalinta
  rolė atsiejama nuo ją pasirinkusių pokalbių. \\
\end{xltabular}

Įmonė priima tokius savo profilio laukus: \code{name}, \code{registrationCode}, \code{country},
\code{legalForm}, \code{naceCodes}, \code{sector}, \code{sizeCategory}, \code{employees},
\code{annualTurnoverEur}, \code{activities}, \code{usesAi}, \code{providesAi},
\code{aiDescription}, \code{processesPersonalData}, \code{isPublicSector},
\code{institutionType}, \code{financialMarket} ir \code{attributes}. Keturi taip ar ne tipo laukai
priima \code{true}, \code{false} arba \code{null}, ir \code{null} reiškia \emph{nežinoma}: API
niekada jų nespėja, o eilutę \code{"false"} atmeta, užuot ją perskaičiusi kaip teigiamą reikšmę.

\subsection{Dokumentai}\label{sec:api:documents}

Įmonės dokumentai: įkėlimas, paieška ir skaitymas (\file{companies/api/document_views.py};
\cref{ch:company}). Dokumentas priklauso jį įkėlusiam naudotojui ir gali būti susietas su įmone
įkeliant ar vėliau.

\begin{xltabular}{\textwidth}{@{}L{19mm}Y@{}}
\caption{Maršrutai: dokumentai}\label{tab:api:documents}\\
\toprule \hd{Kas} & \hd{Ką daro maršrutas} \\ \midrule \endfirsthead
\toprule \hd{Kas} & \hd{Ką daro maršrutas} \\ \midrule \endhead
\bottomrule \endlastfoot
\apiroute{GET POST}{/api/documents}
savininkas & \code{GET}: \code{documents} ir \code{total}, naujausi pirmi, be teksto; filtrai
  \code{companyId}, \code{q} (pavadinimas), \code{scope=own}. \code{POST} yra
  \code{multipart/form-data}: \code{file} ir neprivalomai \code{companyId}, \code{title} (jo
  nesant --- failo vardas), \code{kind} (\code{policy}, \code{procedure}, \code{contract},
  \code{register}, \code{other}). Tekstas išgaunamas prieš atsakymą --- 201 su dokumentu; jo
  \code{extractionStatus} yra \code{done}, \code{empty} (teksto nėra, kaip skenuotame PDF) arba
  \code{failed}. Kai tekstas yra, į eilę įtraukiama užduotis \code{documents.index}.
  \emph{Klaidos:} 400 \code{MISSING_FILE}, \code{INVALID_FIELD}; 413 \code{FILE_TOO_LARGE},
  \code{QUOTA_EXCEEDED}. \\
\addlinespace
\apiroute{GET}{/api/documents/search}
savininkas & Kviečiančiojo dokumentų ištraukos, atitinkančios \code{q} (privalomas, daugiausia
  1\,000 simbolių), kaip \code{hits} su dokumentu, antrašte, tekstu ir \code{charStart}. \code{k}
  --- rezultatų skaičius, nuo 1 iki 30, jo nesant 8; \code{companyId}, \code{documentIds}
  (atskirti kableliais, daugiausia 50) ir \code{scope=own} susiaurina paiešką. \emph{Klaidos:} 400
  \code{INVALID_QUERY}. \\
\addlinespace
\apiroute{GET PATCH DELETE}{/api/documents/{id}}
savininkas & \code{GET}: metaduomenys su \code{outline} (antraštės ir jų poslinkiai) ir
  \code{inventory}. \code{PATCH}: \code{companyId} (\code{null} dokumentą atsieja),
  \code{title}, \code{kind}. \code{DELETE} pašalina eilutę, o paskui failą. \emph{Klaidos:} 400
  \code{EMPTY_UPDATE}, \code{INVALID_FIELD}. \\
\addlinespace
\apiroute{GET}{/api/documents/{id}/text}
savininkas & Išgauto teksto langas: \code{offset} (jo nesant 0) ir \code{limit} (jo nesant
  20\,000, daugiausia 100\,000). Grąžina \code{text}, \code{offset}, \code{limit}, \code{total}
  ir \code{extractionStatus}. \emph{Klaidos:} 400 \code{INVALID_QUERY}. \\
\end{xltabular}

\subsection{Pokalbis}\label{sec:api:chat}

API saugo pokalbius ir jų žinutes (\file{chat/api/views.py}); į patį ėjimą atsako asistento
paslauga (\cref{sec:api:assistant}, \cref{ch:chat}). Pokalbius mato tik jų naudotojas.

\begin{xltabular}{\textwidth}{@{}L{19mm}Y@{}}
\caption{Maršrutai: pokalbiai ir žinutės}\label{tab:api:chat}\\
\toprule \hd{Kas} & \hd{Ką daro maršrutas} \\ \midrule \endfirsthead
\toprule \hd{Kas} & \hd{Ką daro maršrutas} \\ \midrule \endhead
\bottomrule \endlastfoot
\apiroute{GET}{/api/chat/models}
prisijungęs & Modeliai, kuriais pokalbis gali atsakyti: \code{models} --- DI tinklų sąsajos
  pokalbių modeliai su \code{id}, \code{name} ir \code{pricing}, ir \code{default} --- aktyvios
  \prompt{chat.assistant} versijos modelis. Kol DI tinklų sąsajos pasiekti nepavyksta, sąrašas
  tuščias. \\
\addlinespace
\apiroute{GET POST}{/api/chat/threads}
tik savininkas & \code{GET}: kviečiančiojo pokalbiai, kurie nėra archyvuoti, neseniausiai
  atnaujinti pirmi, daugiausia 300; filtras \code{companyId}. \code{POST}: neprivalomi
  \code{companyId} ir \code{title}; atsako 201 su pokalbiu. \\
\addlinespace
\apiroute{GET PATCH DELETE}{/api/chat/threads/{id}}
tik savininkas & \code{PATCH}: \code{title} (daugiausia 200 simbolių), \code{companyId}
  (\code{null} atsieja įmonę; pakeitus įmonę rolė atmetama), \code{roleId} (pokalbio įmonės
  rolė), \code{model} (vienas iš \code{/api/chat/models} identifikatorių arba \code{null}
  numatytajam). \code{DELETE} pokalbį archyvuoja; eilutės lieka. \emph{Klaidos:} 400
  \code{UNKNOWN_MODEL}, \code{EMPTY_UPDATE}; 503 \code{MODELS_UNAVAILABLE}. \\
\addlinespace
\apiroute{GET POST}{/api/chat/threads/{id}/messages}
tik savininkas & \code{GET}: istorija --- paprastas eilučių \code{id}, \code{parent_id},
  \code{format}, \code{content} sąrašas. \code{POST} įrašo arba pakeičia vieną žinutę pagal jos
  identifikatorių pokalbyje ir atsako \code{ok}, \code{created} ir \code{autoTitle} (ar pokalbis
  vis dar turi sau suteiktą pavadinimą). Čia galima rašyti tik paties naudotojo žinutes;
  asistento atsakymai ateina per vidines duris. \emph{Klaidos:} 403 \code{ROLE_NOT_ALLOWED}; 409
  \code{MESSAGE_ID_TAKEN}; 413 \code{MESSAGE_TOO_LARGE}. \\
\end{xltabular}

\subsection{Teisės aktų bazė}\label{sec:api:corpus}

Teisės aktai, jų redakcijos ir nuostatos, paieška ir nuorodų tikrintuvas
(\file{corpus/api/views.py}; \cref{ch:corpus}). Skaityti gali kiekvienas prisijungęs; naudotojas
neįjungtam teisės aktui gauna 404. Rašyti gali tik administratoriai, ir tai įrašoma į veiksmų žurnalą.
Joks maršrutas čia pats nieko neįkelia: teisės aktą apdorojantys maršrutai įtraukia užduotį į eilę
ir atsako 202 su ja (\cref{sec:api:client}).

\begin{xltabular}{\textwidth}{@{}L{19mm}Y@{}}
\caption{Maršrutai: teisės aktų bazė}\label{tab:api:corpus}\\
\toprule \hd{Kas} & \hd{Ką daro maršrutas} \\ \midrule \endfirsthead
\toprule \hd{Kas} & \hd{Ką daro maršrutas} \\ \midrule \endhead
\bottomrule \endlastfoot
\apiroute{GET POST}{/api/corpus/documents}
prisijungęs; \code{POST}: administratorius & \code{GET}: katalogas kaip \code{documents}, kiekvienas
  teisės aktas su dabartine redakcija, būsena, paskutine klaida ir pareigų skaičiumi; filtrai
  \code{q}, \code{jurisdiction}, \code{docType}, \code{domain}. \code{POST} prideda teisės aktą:
  \code{title}, \code{jurisdiction} (\code{EU}, \code{LT}, \code{INTL}), \code{sourceKind}
  (\code{cellar}, \code{etar}, \code{eseimas}, \code{lb}, \code{url}, \code{upload}; jo nesant
  --- \code{upload}) ir, kaip reikia šaltiniui, \code{externalId} arba https \code{sourceUrl};
  taip pat \code{shortTitle}, \code{titleEn}, \code{docType}, \code{language}, \code{domains},
  \code{appliesTo}, \code{applicabilityNotes}, \code{status}, \code{enabled}. Teisės akto, kurio
  šaltinį galima atsisiųsti, įkėlimas iš karto įtraukiamas į eilę; \code{upload} teisės aktas
  laukia savo failo. \emph{Klaidos:} 400 \code{INVALID_FIELD}; 409 \code{SLUG_TAKEN}. \\
\addlinespace
\apiroute{GET PATCH DELETE}{/api/corpus/documents/{id}}
prisijungęs; \code{PATCH}, \code{DELETE}: administratorius & \code{GET}: teisės aktas su
  \code{versions} ir \code{latestJob}, taip pat tai, ką modeliai perskaitė iš paties dabartinės
  redakcijos teksto: \code{scopeReading} (\code{types}, \code{reasoning}, \code{readAt}) ir
  \code{sanctionsReading} (\code{tiers} --- pakopų skaičius, \code{note}, \code{readAt});
  kiekvienas yra \code{null}, kol dar neskaityta. \code{PATCH} redaguoja tuos pačius laukus. \code{DELETE}
  pašalina teisės aktą su kiekviena redakcija, nuostata, fragmentu ir saugomu failu. \\
\addlinespace
\apiroute{POST}{/api/corpus/documents/{id}/refresh}
administratorius & Įtraukia į eilę atsisiuntimą ir įkėlimą. \code{force} tekstą apdoroja iš naujo
  net tada, kai šaltinis nepasikeitė; \code{structure} (\code{auto}, \code{rules}, \code{ai},
  \code{pages}; tik standartui) nurodo, kaip šį kartą atpažįstama standarto struktūra. Atsako
  202 su užduotimi --- jau esančia eilėje ar vykdoma, jei tokia yra. \emph{Klaidos:} 400
  \code{INVALID_FIELD}; 409 \code{NO_FILE} (\code{upload} teisės aktas be failo). \\
\addlinespace
\apiroute{POST}{/api/corpus/documents/{id}/upload}
administratorius & \code{multipart/form-data} su \code{file}: failas tampa nauja redakcija, ir jos
  įkėlimas įtraukiamas į eilę; 202 su užduotimi. \emph{Klaidos:} 400 \code{MISSING_FILE}; 409
  \code{UNCHANGED} (tie patys baitai kaip dabartinės redakcijos); 413 \code{FILE_TOO_LARGE}. \\
\addlinespace
\apiroute{GET}{/api/corpus/documents/{id}/units}
prisijungęs & Dabartinės redakcijos arba redakcijos \code{version} nuostatos teisės akto tvarka,
  su sutrumpintais tekstais: \code{units}, \code{total}, \code{offset}, \code{limit}. Filtrai
  \code{type} (kableliais atskirtas nuostatų tipų sąrašas) ir \code{q}. \\
\addlinespace
\apiroute{GET}{/api/corpus/units/{id}}
prisijungęs & Viena nuostata visa, su jos teisės aktu, \code{versionNo}, \code{isCurrent} ir
  \code{citationUrl} --- nuoroda į oficialų tekstą. \\
\addlinespace
\apiroute{GET}{/api/corpus/provision}
prisijungęs & Tikslus vienos nuostatos tekstas, randamas pagal nuorodą: \code{document} (trumpasis
  pavadinimas, sisteminis identifikatorius, identifikatorius ar oficialus numeris), \code{article} (arba
  \code{clause} --- standartui), \code{paragraph}, \code{type}. Įterptas straipsnis randamas,
  kad ir kaip būtų užrašytas jo numeris (\enquote{12-1}, \enquote{12¹}, \enquote{12(1)}). 404,
  kai teisės aktų bazėje
  tokios nuostatos nėra. \\
\addlinespace
\apiroute{POST}{/api/corpus/search}
prisijungęs & Paieška dabartinėse redakcijose. Turinys \code{query} ir neprivalomai \code{k} (nuo
  1 iki 30, jo nesant 8), \code{documents} (daugiausia 50 teisės aktų, kiekvienas nurodomas
  taip, kaip jį priima bet kuri kita teisės akto paieška, --- identifikatoriumi, sisteminiu
  identifikatoriumi, trumpuoju pavadinimu ar CELEX numeriu, neatsižvelgiant į raidžių dydį;
  pavadinimas, neatitinkantis jokio teisės akto, paiešką susiaurina tiek, kad ji nieko
  neranda), \code{jurisdiction}, \code{kind} (\code{acts} arba \code{standards}). Grąžina
  \code{hits}. \emph{Klaidos:} 400 \code{MISSING_FIELDS}, \code{INVALID_FIELD}. \\
\addlinespace
\apiroute{POST}{/api/corpus/check-citations}
prisijungęs & Nuorodų tikrintuvas. Turinys \code{citations} --- nuo 1 iki 50 elementų iš
  \code{act}, \code{article} ir neprivalomai \code{paragraph}, \code{point},
  \code{claimedHeading}, kiekvienas daugiausia 300 simbolių. Grąžina \code{results}, po vieną
  kiekvienai nuorodai ta pačia tvarka: \code{status} (\code{found}, \code{not_found},
  \code{act_not_in_corpus}), rastą nuostatą, \code{headingMatch} (\code{match}, \code{mismatch},
  \code{unknown}) ir \code{judgedBy} (\code{model}, \code{overlap} arba \code{null}). Kai
  kviečiančiojo biudžetas išnaudotas arba modelio iškvietimas nepavyksta, teiginys vertinamas
  ne modeliu, o pagal žodžių sutapimą; pati patikra niekada nesustoja. \emph{Klaidos:} 400
  \code{INVALID_FIELD}. \\
\addlinespace
\apiroute{POST}{/api/corpus/reembed}
administratorius & Įtraukia į eilę vieno teisės akto (\code{documentId}) arba, jo nenurodžius,
  visų teisės aktų vektorių perskaičiavimą; 202 su užduotimi. Reikalingas tik pasikeitus
  vektorizavimo modeliui. \\
\addlinespace
\apiroute{POST}{/api/corpus/refresh-all}
administratorius & Įtraukia į eilę kiekvieno šaltinio patikrą, ar nėra naujos redakcijos; 202
  su užduotimi. \\
\end{xltabular}

\subsection{Pareigos}\label{sec:api:obligations}

Pareigų registras ir jo peržiūra (\file{corpus/api/obligation_views.py}; \cref{ch:obligations}).
Skaityti gali kiekvienas prisijungęs; pareigą keičia tik administratorius, ir kiekvienas
pakeitimas įrašomas į veiksmų žurnalą.

\begin{xltabular}{\textwidth}{@{}L{19mm}Y@{}}
\caption{Maršrutai: pareigos}\label{tab:api:obligations}\\
\toprule \hd{Kas} & \hd{Ką daro maršrutas} \\ \midrule \endfirsthead
\toprule \hd{Kas} & \hd{Ką daro maršrutas} \\ \midrule \endhead
\bottomrule \endlastfoot
\apiroute{GET POST}{/api/obligations}
prisijungęs; \code{POST}: administratorius & \code{GET}: \code{obligations} ir \code{total}.
  Filtrai \code{document}, \code{status} (\code{proposed}, \code{approved}, \code{rejected};
  viena būsena arba kelios, atskirtos kableliais),
  \code{category}, \code{q} ir \code{verification} --- vertintojo verdiktas: \code{supported},
  \code{partly}, \code{unsupported}, \code{doubt} (paskutiniai du kartu) arba \code{unchecked}.
  \code{order=review} abejotinas pateikia pirmiausia, o tarp jų --- pirmiausia reikšmingiausias.
  \code{POST} įveda pareigą rankiniu būdu: \code{title}, \code{text}, neprivalomai
  \code{documentId} ir toliau išvardyti redaguojami laukai; ji saugoma su šaltiniu
  \code{standard} ir iš karto patvirtinta. \emph{Klaidos:} 400 \code{INVALID_QUERY},
  \code{MISSING_FIELDS}, \code{INVALID_FIELD}. \\
\addlinespace
\apiroute{POST}{/api/obligations/review}
administratorius & Masinis sprendimas: \code{ids} (nuo 1 iki 500) ir \code{status}. Kiekviena
  pareiga pažymima peržiūrėtoju; viena veiksmų žurnalo eilutė jas išvardija. Grąžina
  \code{updated} ir \code{status}. \emph{Klaidos:} 400 \code{INVALID_FIELD}. \\
\addlinespace
\apiroute{GET PATCH DELETE}{/api/obligations/{id}}
prisijungęs; \code{PATCH}, \code{DELETE}: administratorius & \code{PATCH}: bet kuris iš
  \code{title}, \code{titleEn}, \code{text}, \code{textEn}, \code{addressee}, \code{conditions},
  \code{evidence}, \code{deadline}, \code{notes}, \code{category}, \code{severity},
  \code{appliesTo}, \code{status}. Nustačius \code{status}, pažymimas peržiūrėtojas. Pakeistas
  pavadinimas ar tekstas vektorizuojamas iš naujo ir netenka vertintojo verdikto, kuris buvo
  duotas kitiems žodžiams. \emph{Klaidos:} 400 \code{INVALID_FIELD}. \\
\addlinespace
\apiroute{POST}{/api/corpus/documents/{id}/extract-obligations}
administratorius & Įtraukia į eilę pareigų išskyrimą iš dabartinės teisės akto redakcijos;
  \code{language}. 202 su užduotimi arba su ta, kuri jau vykdoma. \emph{Klaidos:} 409
  \code{NOT_INGESTED} (teisės aktas dar neturi redakcijos). \\
\addlinespace
\apiroute{POST}{/api/corpus/documents/{id}/verify-obligations}
administratorius & Įtraukia į eilę vertintojo patikrą tų dabartinės redakcijos pareigų, kurių
  jis dar neįvertino, --- su \code{again: true} visų. 202 su užduotimi arba su ta, kuri jau vykdoma.
  \emph{Klaidos:} 400 \code{INVALID_FIELD}; 409 \code{NOT_INGESTED}. \\
\end{xltabular}

\subsection{Ataskaitos}\label{sec:api:reports}

Atitikties ataskaitos paleidimas, apimties peržiūra, išvadų peržiūra, patvirtinimas ir eksportai
(\file{analysis/api/views.py}; \cref{ch:report}). Ataskaita seka paskui savo įmonę: ją mato įmonės
savininkas ir administratoriai. Vienu metu vienai įmonei rengiama viena ataskaita.

\begin{xltabular}{\textwidth}{@{}L{19mm}Y@{}}
\caption{Maršrutai: ataskaitos}\label{tab:api:reports}\\
\toprule \hd{Kas} & \hd{Ką daro maršrutas} \\ \midrule \endfirsthead
\toprule \hd{Kas} & \hd{Ką daro maršrutas} \\ \midrule \endhead
\bottomrule \endlastfoot
\apiroute{GET POST}{/api/reports}
savininkas & \code{GET}: \code{reports}, naujausios pirmos, daugiausia 200, kiekviena su savo
  užduoties \code{jobStatus} ir \code{progress}; filtrai \code{companyId} ir \code{status}
  (ataskaitos būsena), taip pat \code{scope=own}, kuris administratoriaus sąrašą susiaurina
  iki jo paties įmonių ataskaitų, --- visi pritaikomi prieš sutrumpinant sąrašą iki tų 200. \code{POST} paleidžia
  ataskaitą: \code{companyId}, \code{language} ir neprivalomai \code{documentIds} (daugiausia 200
  įjungtų teisės aktų, kuriais ataskaita apribojama), \code{review} (sustoti apimties peržiūrai)
  ir \code{internal} (vertinti ir įmonės vidaus reikalavimus). Atsako 201 su į eilę įtraukta
  ataskaita; užduotis yra \code{report.analyze}. \emph{Klaidos:} 400 \code{INVALID_FIELD}; 409
  \code{CORPUS_EMPTY}, \code{ALREADY_RUNNING} (ir tada, kai įmonės ataskaita laukia apimties
  patvirtinimo); 429 \code{BUDGET_EXCEEDED}. \\
\addlinespace
\apiroute{GET}{/api/reports/estimate}
savininkas & Kiek dabar kainuotų įmonės \code{companyId} ataskaita --- viršutinė riba JAV
  doleriais, su kviečiančiojo \code{budget}; \code{documentIds} ir \code{internal=true} kaip
  paleidžiant. \emph{Klaidos:} 400 \code{INVALID_FIELD}, kai \code{documentIds} atmestų ir
  paleidimas (teisės akto nėra teisės aktų bazėje arba jis išjungtas); 404 \code{NOT_FOUND}. \\
\addlinespace
\apiroute{GET PATCH DELETE}{/api/reports/{id}}
savininkas & \code{GET}: visa ataskaita --- \code{acts}, \code{findings}, \code{actions},
  \code{citations}, \code{summary}, \code{stats}, patvirtinimai ir užduoties būsena bei eiga.
  \code{PATCH} yra patvirtinimas: \code{approved} (\code{true} arba \code{false}) ir
  \code{note}; atsakymas papildomas \code{register} --- tuo, ką patvirtinimas pakeitė veiksmų
  registre. \code{DELETE} pirmiausia atšaukia patvirtinimą. \emph{Klaidos:} 400
  \code{MISSING_FIELDS}; 409 \code{NOT_FINISHED} (patvirtinimas), \code{STILL_RUNNING}
  (šalinimas). \\
\addlinespace
\apiroute{POST}{/api/reports/{id}/cancel}
savininkas & Sustabdo eilėje esančią ar vykdomą ataskaitą, o apimties patvirtinimo laukiančią
  atšaukia iš karto. \emph{Klaidos:} 409 \code{NOT_CANCELLABLE}. \\
\addlinespace
\apiroute{POST}{/api/reports/{id}/scope}
savininkas & Ataskaitos, laukiančios būsenoje \code{scope_review}, apimties peržiūra. Turinys
  \code{acts} --- žemėlapis iš teisės akto identifikatoriaus ataskaitoje į \code{assess} arba
  \code{skip} (nepaminėtam teisės aktui paliekamas vykdymo sprendimas) --- ir \code{approved}. Be
  \code{approved} nieko neįrašoma, o atsakymas įkainoja pasirinkimą: \code{estimate} ir
  \code{budget}. Su \code{approved: true} pasirinkimai pritaikomi, patvirtinimas pažymimas ir su
  kiekvienu pakeitimu įrašomas į veiksmų žurnalą, o užduotis įtraukiama į eilę antrajai daliai;
  atsakymas yra ataskaita. \emph{Klaidos:} 400 \code{INVALID_FIELD}, \code{NOTHING_TO_ASSESS};
  409 \code{NOT_WAITING}; 429 \code{BUDGET_EXCEEDED}. \\
\addlinespace
\apiroute{PATCH}{/api/reports/{id}/findings/{fid}}
savininkas & Žmogaus sprendimas dėl vienos išvados: \code{overrideStatus} (\code{met},
  \code{partial}, \code{missing}, \code{unknown}, \code{not_applicable} arba \code{null}, kad jį
  panaikintų), \code{reviewerNote} (4\,000 simbolių), \code{cause} (2\,000). Bet kuris iš jų
  atšaukia patvirtinimą; ataskaitos suvestiniai skaičiai perskaičiuojami. Įrašoma į veiksmų
  žurnalą. \emph{Klaidos:} 400 \code{INVALID_FIELD}. \\
\addlinespace
\apiroute{GET}{/api/reports/{id}/export}
savininkas & Užbaigta ataskaita kaip failas: \code{format} yra \code{pdf} (jo nesant) arba
  \code{xlsx}. Įrašoma į veiksmų žurnalą. \emph{Klaidos:} 400 \code{INVALID_QUERY}; 409
  \code{NOT_FINISHED}. \\
\end{xltabular}

\subsection{Veiksmų registras ir teisės pokyčiai}\label{sec:api:actions}

Tai, kas lieka gyventi po ataskaitos patvirtinimo (\file{analysis/api/views.py},
\file{analysis/register.py}; \cref{ch:register}). Pats registras keičiasi tik patvirtinus
ataskaitą; šie maršrutai jį skaito ir leidžia žmogui dirbti su atskirais veiksmais.

\begin{xltabular}{\textwidth}{@{}L{19mm}Y@{}}
\caption{Maršrutai: veiksmų registras ir teisės pokyčiai}\label{tab:api:actions}\\
\toprule \hd{Kas} & \hd{Ką daro maršrutas} \\ \midrule \endfirsthead
\toprule \hd{Kas} & \hd{Ką daro maršrutas} \\ \midrule \endhead
\bottomrule \endlastfoot
\apiroute{GET}{/api/companies/{id}/actions}
savininkas & Įmonės veiksmų registras: \code{actions} (daugiausia 500) --- pirmiausia
  vėluojantys, tada kiti atviri pagal terminą, tada atlikti ir uždaryti; ir \code{counts} pagal
  būseną (\code{late}, \code{open}, \code{in_progress}, \code{done}, \code{closed}). \\
\addlinespace
\apiroute{GET}{/api/companies/{id}/actions/export}
savininkas & Registras kaip darbaknygė (\code{.xlsx}); \code{language}. Įrašoma į veiksmų
  žurnalą. \\
\addlinespace
\apiroute{PATCH}{/api/actions/{id}}
savininkas & Žmogaus darbas su vienu veiksmu: \code{status} (\code{open}, \code{in_progress},
  \code{done}, \code{closed}), \code{owner}, \code{dueDate}, \code{note}. Įrašoma į veiksmų
  žurnalą. \emph{Klaidos:} 400 \code{INVALID_FIELD}, \code{NOTE_REQUIRED} (uždarant rankiniu būdu
  be pastabos); 409 \code{CLOSED_BY_REPORT} (iš naujo atidarant veiksmą, kurį uždarė patvirtinta
  ataskaita). \\
\addlinespace
\apiroute{GET}{/api/companies/{id}/law-changes}
savininkas & Pagal naujausią užbaigtą įmonės ataskaitą: vertinti teisės aktai, kurie nuo tada
  turi naujesnę redakciją, jų pasikeitusios nuostatos ir jomis paremtos išvados. Grąžina
  \code{reportId}, \code{reportDate}, \code{acts}. \\
\end{xltabular}

\subsection{Atitikties programa}\label{sec:api:programme}

Įmonės programos dokumentai --- rūšys \code{policy}, \code{governance} ir \code{risks}
(\file{programme/api/views.py}; \cref{ch:programme}). Nuosavybė seka paskui įmonę. Vienai įmonei
vienu metu rengiamas vienas tos pačios rūšies dokumentas.

\begin{xltabular}{\textwidth}{@{}L{19mm}Y@{}}
\caption{Maršrutai: programos dokumentai}\label{tab:api:programme}\\
\toprule \hd{Kas} & \hd{Ką daro maršrutas} \\ \midrule \endfirsthead
\toprule \hd{Kas} & \hd{Ką daro maršrutas} \\ \midrule \endhead
\bottomrule \endlastfoot
\apiroute{GET POST}{/api/companies/{id}/programme}
savininkas & \code{GET}: \code{products} --- pagal rūšį, naujausia versija pirma, be turinio,
  kiekvienas su savo užduoties būsena ir eiga; ir \code{kinds} --- trys rūšys su pavadinimais ir
  tuo, kuo kiekviena remiasi. \code{POST}: \code{kind} ir \code{language}; paleidžia naujos
  versijos rengimą (užduotis \code{programme.generate}) ir atsako 201 su į eilę įtrauktu
  dokumentu. \emph{Klaidos:} 400 \code{INVALID_FIELD}; 409 \code{NO_DOCUMENTS} (įmonė neturi
  perskaitomo dokumento), \code{ALREADY_RUNNING}; 429 \code{BUDGET_EXCEEDED}. \\
\addlinespace
\apiroute{GET PATCH DELETE}{/api/programme/{id}}
savininkas & \code{GET}: dokumentas su \code{content}. \code{PATCH} su \code{approved: true} ir
  neprivalomai \code{note} yra žmogaus atliekamas projekto patvirtinimas. \code{DELETE} pašalina
  versiją. \emph{Klaidos:} 400 \code{MISSING_FIELDS}; 409 \code{NOT_DRAFT}, \code{RULES_BROKEN}
  (pažeista blokuojanti tos rūšies taisyklė; sakinys įvardija taisykles), \code{STILL_RUNNING}
  (šalinimas). \\
\addlinespace
\apiroute{PUT}{/api/programme/{id}/content}
savininkas & Žmogaus atliktas projekto pakeitimas, siunčiamas visas kaip \code{content}. Turinys
  tikrinamas pagal savo rūšies taisykles be modelio iškvietimo, o atsakymas yra dokumentas taip,
  kaip įrašytas, su iš naujo apskaičiuotomis išvestinėmis reikšmėmis. \emph{Klaidos:} 400
  \code{MISSING_FIELDS}, \code{INVALID_FIELD}; 409 \code{NOT_DRAFT}, \code{UNKNOWN_KIND}. \\
\addlinespace
\apiroute{POST}{/api/programme/{id}/revise}
savininkas & Naujas projektas, viena versija aukštesnis, nukopijuotas iš patvirtinto dokumento,
  kuris lieka patvirtintas; 201. \emph{Klaidos:} 409 \code{NOT_APPROVED}, \code{NEWER_VERSION}. \\
\addlinespace
\apiroute{POST}{/api/programme/{id}/cancel}
savininkas & Sustabdo eilėje esantį ar vykdomą rengimą. \emph{Klaidos:} 409
  \code{NOT_CANCELLABLE}. \\
\addlinespace
\apiroute{GET}{/api/programme/{id}/chart}
savininkas & Valdymo struktūros dokumento struktūra, nubraižyta kaip organizacinė schema,
  \code{image/svg+xml}. 404, kai dokumentas neturi struktūros. \\
\addlinespace
\apiroute{GET}{/api/programme/{id}/export}
savininkas & Dokumentas kaip failas: \code{format} yra \code{pdf} bet kuriai rūšiai, taip pat
  \code{docx} politikai ir \code{xlsx} rizikos registrui. Projektas eksportuojamas taip pat, kaip
  ir patvirtintas dokumentas. Įrašoma į veiksmų žurnalą. \emph{Klaidos:} 400
  \code{INVALID_QUERY}; 409 \code{NOT_READY}. \\
\end{xltabular}

\subsection{Instrukcijos}\label{sec:api:prompts}

Instrukcijų ir jų versijų saugykla administratoriams (\file{prompts/api/views.py};
\cref{ch:prompts}). Kiekvienas įrašymas ir kiekvienas bandomasis paleidimas įrašomas į veiksmų
žurnalą. Modelių sąrašas, iš kurio pildomas redaktoriaus modelio pasirinkimas, yra faile
\file{assistant/api/mcp_views.py}.

\begin{xltabular}{\textwidth}{@{}L{19mm}Y@{}}
\caption{Maršrutai: instrukcijos ir modeliai}\label{tab:api:prompts}\\
\toprule \hd{Kas} & \hd{Ką daro maršrutas} \\ \midrule \endfirsthead
\toprule \hd{Kas} & \hd{Ką daro maršrutas} \\ \midrule \endhead
\bottomrule \endlastfoot
\apiroute{GET}{/api/prompts}
administratorius & Visos instrukcijos kaip \code{prompts}: \code{slug}, \code{name}, \code{kind},
  \code{variables}, \code{activeVersion}, \code{activeModel}, \code{pendingDefault}. \\
\addlinespace
\apiroute{GET}{/api/prompts/{slug}}
administratorius & Viena instrukcija su \code{versions}, naujausios pirmos, kiekviena su savo
  tekstu. \\
\addlinespace
\apiroute{POST}{/api/prompts/{slug}/versions}
administratorius & Prideda versiją: \code{text} (daugiausia 60\,000 simbolių), \code{model} ir
  neprivalomai \code{temperature} (nuo 0 iki 2), \code{maxTokens} (nuo 1 iki 128\,000),
  \code{note} (500 simbolių), \code{activate}. Atsako 201 su versija. \emph{Klaidos:} 400
  \code{MISSING_FIELDS}, \code{INVALID_FIELD}. \\
\addlinespace
\apiroute{POST}{/api/prompts/{slug}/versions/{n}/activate}
administratorius & Padaro versiją \code{n} aktyvia; grąžina instrukciją. \\
\addlinespace
\apiroute{POST}{/api/prompts/{slug}/test}
administratorius & Vienas modelio iškvietimas, kai instrukcija yra įrašyta versija
  (\code{versionNo}; jo nesant --- aktyvioji) arba neįrašytas projektas (\code{text},
  \code{model}), o duomenys --- \code{input} (daugiausia 200\,000 simbolių); naudojama
  instrukcijos etapo išvesties schema, jei instrukcija ją turi; \code{variables} užpildo
  kintamuosius. Grąžina \code{text}, \code{object}, \code{model}, \code{finishReason},
  \code{usage}, \code{costUsd}, \code{durationMs}, \code{renderedInstructions}. Atsakęs
  iškvietimas priskaičiuojamas prie kviečiančiojo biudžeto ir įrašomas kaip modelio
  iškvietimas; nepavykęs tik įrašomas į veiksmų žurnalą. \emph{Klaidos:} 400
  \code{MISSING_FIELDS}, \code{INVALID_FIELD}; 429 \code{BUDGET_EXCEEDED}; 502
  \code{MODEL_ERROR}. \\
\addlinespace
\apiroute{GET}{/api/llm/models}
administratorius & DI tinklų sąsajos siūlomi modeliai taip, kaip juos pateikia asistento
  paslauga: \code{providers} ir \code{models}, kiekvienas su \code{id}, \code{name},
  \code{kind} (\code{chat} arba \code{embedding}), \code{contextLength}, \code{maxOutputTokens}
  ir \code{pricing}. \\
\end{xltabular}

\subsection{MCP serveriai}\label{sec:api:mcp}

Išorinių įrankių serverių, kurių įrankius gali kviesti pokalbių asistentas, registras
administratoriams (\file{assistant/api/mcp_views.py}; \cref{ch:chat}). Serverio antraščių
reikšmės yra prisijungimo duomenys: jos niekada negrąžinamos, grąžinami tik jų pavadinimai
(\code{headerNames}).

\begin{xltabular}{\textwidth}{@{}L{19mm}Y@{}}
\caption{Maršrutai: MCP serveriai}\label{tab:api:mcp}\\
\toprule \hd{Kas} & \hd{Ką daro maršrutas} \\ \midrule \endfirsthead
\toprule \hd{Kas} & \hd{Ką daro maršrutas} \\ \midrule \endhead
\bottomrule \endlastfoot
\apiroute{GET POST}{/api/mcp/servers}
administratorius & \code{POST}: \code{name} (nuo 2 iki 31 simbolio: mažosios raidės, skaitmenys ir
  pabraukimo brūkšnys, pirma --- raidė), \code{url} ir neprivalomai \code{transport} (\code{http}
  arba \code{sse}), \code{headers} (daugiausia 10), \code{allowedTools} (tuščias reiškia visus
  serverio įrankius), \code{description}, \code{enabled}. URL turi būti https, standartiniu
  prievadu, į viešą DNS vardą --- jokio IP adreso, jokio konteinerio ar vietinio vardo, jokių
  prisijungimo duomenų jame. Naujas serveris iš pradžių būna išjungtas. \emph{Klaidos:} 400
  \code{MISSING_FIELDS}, \code{INVALID_FIELD}; 409 \code{DUPLICATE_NAME}. \\
\addlinespace
\apiroute{GET PATCH DELETE}{/api/mcp/servers/{id}}
administratorius & \code{PATCH}: bet koks tų pačių laukų poaibis. Antraštės sujungiamos:
  reikšmė \code{null} tą antraštę pašalina, kitos išlaikomos. Serverio, kurio įrašytas URL
  pažeidžia taisyklę, įjungti negalima. \emph{Klaidos:} 400 \code{INVALID_FIELD}; 409
  \code{DUPLICATE_NAME}. \\
\addlinespace
\apiroute{POST}{/api/mcp/servers/{id}/test}
administratorius & Asistento paslauga prisijungia prie serverio ir išvardija jo įrankius.
  Atsakymas visada yra 200 su serveriu ir \code{ok}: pavykus \code{tools} pateikia sąrašą,
  nepavykus \code{lastError} nurodo priežastį. Bet kuriuo atveju įrašoma į veiksmų žurnalą.
  \emph{Klaidos:} 400 \code{INVALID_FIELD} (įrašytas URL pažeidžia taisyklę). \\
\end{xltabular}

\subsection{Užduotys}\label{sec:api:jobs}

Užduočių vykdytojo eilės stebėjimas (\file{jobs/api/views.py}; \cref{ch:arch}). Naudotojas mato
savo sukeltas užduotis, administratorius --- visas.

\begin{xltabular}{\textwidth}{@{}L{19mm}Y@{}}
\caption{Maršrutai: užduotys}\label{tab:api:jobs}\\
\toprule \hd{Kas} & \hd{Ką daro maršrutas} \\ \midrule \endfirsthead
\toprule \hd{Kas} & \hd{Ką daro maršrutas} \\ \midrule \endhead
\bottomrule \endlastfoot
\apiroute{GET}{/api/jobs}
savininkas & \code{jobs}, naujausios pirmos; filtrai \code{status}, \code{kind}, \code{subject};
  \code{limit}. Su teisės aktu susijusi užduotis turi teisės akto pavadinimą lauke \code{subjectTitle}.
  \emph{Klaidos:} 400 \code{INVALID_QUERY}. \\
\addlinespace
\apiroute{GET}{/api/jobs/counts}
savininkas & Kiek kviečiančiajam matomų užduočių yra kiekvienoje būsenoje: \code{queued},
  \code{running}, \code{done}, \code{failed}, \code{cancelled}. \\
\addlinespace
\apiroute{GET}{/api/jobs/{id}}
savininkas & Viena užduotis su \code{steps}: kiekvieno etapo pavadinimas, būsena, trukmė, klaida,
  instrukcijos versija, su kuria jis vykdytas, ir ką sunaudojo jo modelių iškvietimai
  (\code{usage}). \\
\addlinespace
\apiroute{POST}{/api/jobs/{id}/cancel}
savininkas & Eilėje esanti užduotis atšaukiama iš karto; vykdoma gauna \code{cancelRequested} ir
  sustoja per kitą savo patikrą. Įrašoma į veiksmų žurnalą. \emph{Klaidos:} 409
  \code{NOT_CANCELLABLE} (jau baigta). \\
\end{xltabular}

\subsection{DI sąnaudos}\label{sec:api:usage}

Kiek kainavo modelių iškvietimai, pagal dieną ir pagal žmogų (\file{assistant/api/usage_views.py};
\cref{ch:models}). Kiekvienas skaito savo sąnaudas. Administratorius praplečia vaizdą iki visų
paskyrų su \code{scope=all} arba žiūri vieną paskyrą su \code{userId}; naudotojas, paprašęs bet
kurio iš jų, gauna 403 \code{FORBIDDEN}.

\begin{xltabular}{\textwidth}{@{}L{19mm}Y@{}}
\caption{Maršrutai: DI sąnaudos}\label{tab:api:usage}\\
\toprule \hd{Kas} & \hd{Ką daro maršrutas} \\ \midrule \endfirsthead
\toprule \hd{Kas} & \hd{Ką daro maršrutas} \\ \midrule \endhead
\bottomrule \endlastfoot
\apiroute{GET}{/api/usage}
savo; administratorius: plačiau & Paskutinės \code{days} dienos (nuo 1 iki 366, jo nesant 30)
  pagal UTC, po eilutę kiekvienai dienai: \code{days}, \code{byKind}, \code{byModel},
  \code{totals}, \code{from}, \code{to}; \code{byUser} administratoriaus plačiajame vaizde; ir
  \code{budget} --- žiūrimos paskyros šiandienos \code{limit}, \code{used} ir \code{remaining}
  (\code{remaining} yra \code{null}, kai biudžetas neribotas). \emph{Klaidos:} 400
  \code{INVALID_QUERY}. \\
\addlinespace
\apiroute{GET}{/api/usage/calls}
savo; administratorius: plačiau & Modelių iškvietimų žurnalas, naujausi pirmi: \code{calls} ir
  \code{total}; filtras \code{kind} (\code{chat}, \code{job_step}, \code{embed},
  \code{prompt_test}). Kiekvienas iškvietimas turi savo modelį, žetonus (įvesties, nuskaitytus
  iš podėlio, įrašytus į jį, išvesties, samprotavimo), \code{costUsd}, trukmę ir baigtį.
  \emph{Klaidos:} 400 \code{INVALID_QUERY}. \\
\addlinespace
\apiroute{GET}{/api/usage/calls/{id}}
savo; administratorius: visi & Vienas iškvietimas su \code{toolCalls} ir, pokalbio ėjimo atveju,
  \code{transcript}: ko naudotojas paklausė ir ką asistentas atsakė. \\
\end{xltabular}

\subsection{Veiksmų žurnalas}\label{sec:api:audit}

Įrašai apie tai, kas ką padarė (\file{audit/api/views.py}; \cref{ch:security}).

\begin{xltabular}{\textwidth}{@{}L{19mm}Y@{}}
\caption{Maršrutai: veiksmų žurnalas}\label{tab:api:audit}\\
\toprule \hd{Kas} & \hd{Ką daro maršrutas} \\ \midrule \endfirsthead
\toprule \hd{Kas} & \hd{Ką daro maršrutas} \\ \midrule \endhead
\bottomrule \endlastfoot
\apiroute{GET}{/api/audit}
administratorius & \code{entries}, naujausi pirmi, kiekvienas su \code{id}, \code{userId},
  \code{userName}, \code{action}, \code{targetType}, \code{targetId}, \code{details}, \code{ip},
  \code{createdAt}; ir \code{next}. Filtrai: \code{action} --- priešdėlis (\code{users.} visiems
  paskyrų pakeitimams) ir \code{user} --- vienos paskyros identifikatorius. \code{next} turi
  \code{before} ir \code{beforeId} kitam puslapiui; juos reikia grąžinti tokius, kokie yra, o po
  paskutinio puslapio jis yra \code{null}. \emph{Klaidos:} 400 \code{INVALID_QUERY}. \\
\end{xltabular}


% ============================================================
\section{Vieši asistento paslaugos maršrutai}\label{sec:api:assistant}
% ============================================================

Asistento paslauga į išorę atsako trimis maršrutais, visais po \code{/api/assistant/chat}
(asistento paslaugoje --- \file{routes/chat.js}; prijungti faile \file{app.js}). Jie priima tą
patį prieigos raktą kaip API: paslauga nustato jo naudotoją, kviesdama \route{GET}{/api/auth/me}, ir
atsakymą prisimena 15 sekundžių. Kas vyksta ėjimo viduje --- instrukcija, įrankiai, modelis, ---
aptariama \cref{ch:chat}.

\begin{xltabular}{\textwidth}{@{}L{19mm}Y@{}}
\caption{Maršrutai: asistento paslauga, vieši}\label{tab:api:assistant}\\
\toprule \hd{Kas} & \hd{Ką daro maršrutas} \\ \midrule \endfirsthead
\toprule \hd{Kas} & \hd{Ką daro maršrutas} \\ \midrule \endhead
\bottomrule \endlastfoot
\apiroute{POST}{/api/assistant/chat}
prisijungęs, savas pokalbis & Vienas pokalbio ėjimas. Turinyje yra viena nauja naudotojo žinutė;
  ankstesnis pokalbis nuskaitomas iš API. Atsakymas yra srautas (\code{text/event-stream})
  pokalbio vykdymo aplinkos naudotojo sąsajos žinučių formatu. \\
\addlinespace
\apiroute{POST}{/api/assistant/chat/{threadId}/stop}
prisijungęs, savas pokalbis & Sustabdo kviečiančiojo vykstantį ėjimą tame pokalbyje. Grąžina
  \code{stopped} --- \code{true} arba \code{false}; dalinis atsakymas įrašomas su būsena
  \code{stopped}. Kito naudotojo pokalbis traktuojamas taip, lyg nebūtų ko stabdyti. \\
\addlinespace
\apiroute{GET}{/api/assistant/chat/{threadId}/status}
prisijungęs, savas pokalbis & Ar pokalbyje vyksta ėjimas: \code{running} ir \code{startedAt}.
  Konsolė jo periodiškai klausia, kai puslapis įkeliamas iš naujo atsakymo viduryje. Niekada
  nelaikomas podėlyje. \\
\end{xltabular}

\paragraph{Ėjimo turinys.} \code{readTurnBody} priima tokį pavidalą ir atkuria jį laukas po
lauko, kad nieko netikėto nepasiektų nei API, nei modelio:

\begin{codeblock}
{"id": "<thread id>",
 "trigger": "submit-message" | "regenerate-message",
 "message": {"id": "<message id>", "role": "user",
             "parts": [{"type": "text", "text": "..."},
                       {"type": "data-document", "data": {"id": "<document id>", "title": "...",
                        "kind": "...", "charCount": 1234, "extractionStatus": "done"}}],
             "metadata": {...}},
 "parentId": "<message id>" | null,
 "locale": "lt" | "en"}
\end{codeblock}

Pokalbio identifikatorius turi būti uuid; \code{trigger} jo nesant yra \code{submit-message};
neprivalomas žinutės \code{metadata} išlaikomas tik tol, kol jis yra objektas, JSON pavidalu
neviršijantis 2\,000 simbolių. Žinutę sudaro nuo 1 iki 30 dalių (\code{MAX_PARTS}); tekstinės
dalys kartu turi daugiausia \code{MAX_MESSAGE_CHARS}~=~60\,000 simbolių. \code{data-document}
dalis yra žymeklis dokumentui, iš anksto įkeltam per \route{POST}{/api/documents}; pats failas
niekada nėra žinutės dalis. Žinutės identifikatorių pasirenka klientas iš raidžių, skaitmenų ir
\code{_ . : -}, daugiausia 128 simbolių.

\paragraph{Prieš srautą.} Pirmiausia veikia turinio analizatorius: 400 \code{INVALID_JSON}
turiniui, kuris nėra JSON, 413 \code{PAYLOAD_TOO_LARGE} viršijus 512~kB. Tada ėjimas nustatyta
tvarka praeina savo patikras, o nepavykus bet kuriai iš jų grąžinama įprasta JSON klaida asistento
paslaugos apvalkale: prieigos raktas (401 \code{AUTH_REQUIRED}, kai jo nėra,
\code{SESSION_INVALID}, kai jis nebegalioja); sukonfigūruotas modelis (503
\code{NOT_CONFIGURED}); užklausų dažnio riba (429 \code{RATE_LIMITED}); turinys (400
\code{INVALID_BODY}, \code{INVALID_THREAD}, \code{INVALID_TRIGGER}, \code{INVALID_MESSAGE},
\code{INVALID_PARENT}, \code{INVALID_PART}, \code{INPUT_TOO_LARGE}, \code{EMPTY_MESSAGE}); dienos
biudžetas (429 \code{BUDGET_EXCEEDED}; jei pati biudžeto patikra nepavyksta, ėjimas vis tiek leidžiamas);
aktyvi \prompt{chat.assistant} versija (503 \code{PROMPT_NOT_CONFIGURED} arba
\code{PROMPT_UNAVAILABLE}, kai jos nepavyksta įkelti); ir pokalbis, gaunamas iš API naudotojo
vardu (404 \code{THREAD_NOT_FOUND}). Kai API pasiekti nepavyksta, atsakymas yra 502
\code{DJANGO_UNREACHABLE}; bet koks kitas API atsisakymas perduodamas su paties API būsena ir
kodu. Išsijungianti paslauga atsako 503 \code{SHUTTING_DOWN}.

\paragraph{Srauto viduje.} Srautui prasidėjus, HTTP būsena yra 200, kad ir kas nutiktų, o
nesėkmė ateina kaip srauto klaidos fragmentas su kodu: \code{MODEL_TIMEOUT} (ėjimas trunka
daugiausia \code{CHAT_TIMEOUT_MS}~=~900 sekundžių), \code{AI_BUDGET_EXHAUSTED} (išnaudotas pačios
DI tinklų sąsajos biudžetas), \code{MODEL_UNAVAILABLE} (laikina nesėkmė: vėliau bandyti dar kartą
gali padėti), \code{MODEL_REFUSED}, \code{MODEL_CONTEXT_EXCEEDED}, \code{MODEL_CONTENT_FILTERED},
\code{MODEL_ACCESS_DENIED} arba \code{STREAM_FAILED} visais kitais atvejais. Ėjimas nepriklauso
nuo ryšio: jei naršyklė atsijungia, atsakymas vis tiek užbaigiamas ir įrašomas, o atsakymo
metaduomenys užfiksuoja, kaip ėjimas baigėsi, jei jis baigėsi nesėkmingai (\code{stopped},
\code{interrupted}, \code{timeout}, \code{budget}, \code{error} arba baigiamosios eilutės, kurią
paslauga turėjo pridėti, rūšis: \code{step-limit}, \code{length}, \code{content-filter},
\code{empty}).


% ============================================================
\section{Vidinės durys}\label{sec:api:internal}
% ============================================================

API ir asistento paslauga reikalingos viena kitai: užduočių vykdytojas ir API kiekvienam modelio
iškvietimui kreipiasi į asistento paslaugą, o asistento paslauga prašo API pokalbio instrukcijos
ir biudžeto ir perduoda jai atsakymus. Šie kreipiniai eina per dvejas \term{vidines duris} ---
kiekvienos paslaugos kelius po \code{/internal/} (\cref{fig:api:doors}). Nė vienų durų negalima
pasiekti iš išorės, ir tam yra dvi nepriklausomos priežastys:

\begin{itemize}
  \item \emph{Ten neveda joks kelias.} Įėjimas persiunčia \code{/api/assistant/*} į asistento
    paslaugą ir \code{/api/*} į API, ir nieko kito nė vienai iš jų. Užklausa
    \code{/internal/...} iš išorės patenka į trečiąją šeimą, ir į ją atsako konsolės
    žiniatinklio serveris. Nė viena iš dviejų paslaugų neskelbia savo prievado.
  \item \emph{Bendra paslaptis.} Kiekvienas vidinis maršrutas reikalauja antraštės
    \code{X-Internal-Secret}, kuri turi sutapti su reikšme, abiem paslaugoms perduodama kaip
    \code{ASSISTANT_INTERNAL_SECRET} paleidžiant diegimą. Palyginimas trunka pastovų laiką, o
    paslauga, kurios paslaptis tuščia, atmeta kiekvieną vidinę užklausą: sutrikus durys
    lieka uždarytos. Atsisakymas yra 403 \code{INTERNAL_ONLY} (\code{require_internal} faile
    \file{assistant/api/internal_views.py}; asistento paslaugoje --- \code{requireInternal}
    faile \file{middleware/internal.js}, kuris veikia dar prieš analizuojant turinį).
\end{itemize}

\begin{figure}[H]
\centering
\begin{tikzpicture}[node distance=8mm and 50mm]
  \node[svc, minimum width=34mm, minimum height=22mm] (django) {API ir\\užduočių\\vykdytojas\\(Django)};
  \node[svc, minimum width=34mm, minimum height=22mm, right=of django] (assistant) {asistento\\paslauga};
  \node[ext, right=14mm of assistant] (gw) {DI tinklų\\sąsaja};
  \draw[call] ([yshift=7mm]django.east) -- ([yshift=7mm]assistant.west)
    node[lbl, midway, above=1pt] {\code{/internal/llm/*}, \code{/internal/mcp/*}};
  \draw[call] (assistant.west) -- (django.east)
    node[lbl, midway] {\code{/internal/assistant/*}};
  \draw[flow] ([yshift=-7mm]assistant.west) -- ([yshift=-7mm]django.east)
    node[lbl, midway, below=1pt] {\code{/api/*} su naudotojo prieigos raktu};
  \draw[call] (assistant) -- (gw);
  \begin{scope}[on background layer]
    \node[lane, fit=(django)(assistant), inner ysep=9mm] (lane) {};
  \end{scope}
  \node[lanelabel, anchor=south west] at (lane.north west) {izoliuotas tinklas};
  \node[note, anchor=north west] at ([yshift=-1mm]lane.south west) {tarp dviejų paslaugų --- punktyrinės: būtina antraštė \code{X-Internal-Secret}; ištisinės: viešoji API, su paties naudotojo teisėmis};
\end{tikzpicture}
\caption{Vidinės durys. API ir jos užduočių vykdytojas kiekvienam modelio iškvietimui kviečia
asistento paslaugą; asistento paslauga kreipiasi atgal dėl to, ką žino tik API. Visa, kas susiję
su paties naudotojo duomenimis, asistento paslauga gauna per viešus maršrutus su to naudotojo
prieigos raktu, todėl nuosavybės taisyklė taikoma vienoje vietoje.}\label{fig:api:doors}
\end{figure}

\subsection{API vidinės durys}\label{sec:api:door-django}

Šeši maršrutai po \code{/internal/assistant/} aptarnauja asistento paslaugą
(\file{assistant/api/internal_views.py}). Jų laukų pavadinimai rašomi snake\_case.

\begin{xltabular}{\textwidth}{@{}L{19mm}Y@{}}
\caption{Maršrutai: API vidinės durys}\label{tab:api:door-django}\\
\toprule \hd{Kas} & \hd{Ką daro maršrutas} \\ \midrule \endfirsthead
\toprule \hd{Kas} & \hd{Ką daro maršrutas} \\ \midrule \endhead
\bottomrule \endlastfoot
\apiroute{GET}{/internal/assistant/prompt}
paslaptis & Instrukcijos, nurodytos \code{slug}, aktyvi versija su tekstu, modeliu ir
  parametrais. 404 \code{PROMPT_NOT_CONFIGURED}, kai instrukcija neturi aktyvios versijos ---
  tada pokalbis atsako 503, užuot veikęs pagal instrukcijas, kurių niekas nemato. \\
\addlinespace
\apiroute{POST}{/internal/assistant/turn-log}
paslaptis & Viena DI sąnaudų žurnalo eilutė pokalbio ėjimui: \code{user_id}, \code{thread_id},
  \code{model}, \code{prompt_version_id}, žetonų skaičiai (\code{input_tokens},
  \code{cached_input_tokens}, \code{cache_write_tokens}, \code{output_tokens},
  \code{reasoning_tokens}), \code{duration_ms}, \code{tool_calls}, \code{outcome} ir
  \code{transcript} (naudotojo tekstas, sutrumpintas iki 20\,000 simbolių, atsakymas --- iki
  100\,000). Kainą apskaičiuoja API. Atsako 201 su \code{id} ir naudotojo \code{budget}. \\
\addlinespace
\apiroute{GET}{/internal/assistant/budget}
paslaptis & Naudotojo \code{user_id} šiandienos biudžetas: \code{limit}, \code{used},
  \code{remaining} (\code{null}, kai neribotas), \code{unit}. Klausiama prieš kiekvieną ėjimą. \\
\addlinespace
\apiroute{POST}{/internal/assistant/threads/{id}/messages}
paslaptis & Įrašo vieną asistento atsakymą: \code{user_id}, \code{id}, \code{parent_id},
  \code{format}, \code{content}. Tai vienintelis asistento žinučių rašytojas, todėl naršyklė
  niekada negali suklastoti įrankio išvesties. Pokalbis turi priklausyti \code{user_id} (kitaip
  404). 201, kai žinutė nauja, 200, kai ji pakeičia esamą. \\
\addlinespace
\apiroute{POST}{/internal/assistant/threads/{id}/title}
paslaptis & Modelio pokalbiui parašytas pavadinimas: \code{user_id}, \code{title}. Jis
  pritaikomas tik tol, kol pokalbis dar turi automatinį pavadinimą, todėl naudotojo nustatytas
  pavadinimas išlieka. Grąžina \code{ok} ir \code{applied}. \\
\addlinespace
\apiroute{GET}{/internal/assistant/mcp-servers}
paslaptis & Įjungti MCP serveriai \emph{su} antraščių reikšmėmis --- vienintelė vieta, kur šie
  prisijungimo duomenys palieka API. Serveris, kurio įrašytas URL pažeidžia adreso taisyklę,
  praleidžiamas. \\
\end{xltabular}

\subsection{Asistento paslaugos vidinės durys}\label{sec:api:door-assistant}

Keturi maršrutai po \code{/internal/} aptarnauja API ir jos užduočių vykdytoją (asistento
paslaugoje --- \file{routes/internal.js} ir \file{routes/internalMcp.js}); klientas API pusėje
yra \file{assistant/gateway.py}. Tai vienintelis sistemos kelias į modelį, ir \cref{ch:models}
juos aprašo išsamiai. Penktajam maršrutui, \code{/health}, paslapties nereikia, ir įėjimas jo taip
pat nepersiunčia.

\begin{xltabular}{\textwidth}{@{}L{19mm}Y@{}}
\caption{Maršrutai: asistento paslaugos vidinės durys}\label{tab:api:door-assistant}\\
\toprule \hd{Kas} & \hd{Ką daro maršrutas} \\ \midrule \endfirsthead
\toprule \hd{Kas} & \hd{Ką daro maršrutas} \\ \midrule \endhead
\bottomrule \endlastfoot
\apiroute{POST}{/internal/llm/generate}
paslaptis & Vienas modelio iškvietimas. Turinys \code{prompt} ir neprivalomai \code{system},
  \code{shared}, \code{model}, \code{schema}, \code{schemaName}, \code{temperature},
  \code{maxOutputTokens}. Grąžina \code{model}, \code{text} arba \code{object} (su schema ---
  pagal ją patikrintas atsakymas), \code{finishReason}, \code{usage}, \code{durationMs}.
  Iškvietimo atsisakoma po \code{GENERATE_TIMEOUT_MS} (numatytai 900 sekundžių).
  \emph{Klaidos:} 400 netinkamai sudarytai užklausai; 402 \code{AI_BUDGET_EXHAUSTED}, kai
  išnaudotas pačios DI tinklų sąsajos biudžetas; kitais atvejais \code{GENERATION_FAILED} su DI
  tinklų sąsajos būsena arba 502 --- ir 429 su \code{retryAfterS}, kai pasiekta užklausų dažnio
  riba. \\
\addlinespace
\apiroute{POST}{/internal/llm/embed}
paslaptis & Vektoriai tekstams \code{texts}, daugiausia 256 per iškvietimą; neprivalomai
  \code{model}. Grąžina \code{model}, \code{embeddings}, \code{dimensions}, \code{usage},
  \code{durationMs}. \emph{Klaidos:} 400 \code{INVALID_TEXTS}, \code{TOO_MANY_TEXTS},
  \code{UNKNOWN_MODEL}; \code{EMBEDDING_FAILED}; 402 \code{AI_BUDGET_EXHAUSTED}. \\
\addlinespace
\apiroute{GET}{/internal/llm/models}
paslaptis & \code{providers} ir \code{models}: DI tinklų sąsajos modelių sąrašas su kainomis,
  naujausi pirmi, laikomas dešimt minučių. \\
\addlinespace
\apiroute{POST}{/internal/mcp/list-tools}
paslaptis & Prisijungia prie MCP serverio, aprašyto \code{url}, \code{transport} ir
  \code{headers}, ir išvardija jo įrankius --- tai \route{POST}{/api/mcp/servers/{id}/test} ryšio
  patikra. \emph{Klaidos:} 400 \code{INVALID_URL}, \code{INVALID_TRANSPORT},
  \code{INVALID_HEADERS}; 502 \code{MCP_UNREACHABLE}, iš pranešimo pašalinus prisijungimo
  duomenis. \\
\addlinespace
\apiroute{GET}{/health}
tinklo viduje & Gyvybingumo ir konfigūracijos patikra (Compose failas būklės patikros nenustato): \code{status},
  \code{model} (\code{configured} arba \code{missing-key}), \code{providers},
  \code{activeTurns}. \\
\end{xltabular}


% ============================================================
\section{Ką turi žinoti klientas}\label{sec:api:client}
% ============================================================

Du API elgesio bruožai nematomi viename maršrute: darbas, trunkantis ilgiau nei jį pradėjusi
užklausa, ir atsakymai, kurie yra failai.

\subsection{Ilgas darbas: paleisti, tikrinti, nuspręsti, atšaukti}\label{sec:api:polling}

Joks maršrutas neverčia kliento laukti įkėlimo, išskyrimo, ataskaitos ar programos dokumento.
Maršrutas užfiksuoja, ko norima, įtraukia užduotį į užduočių vykdytojo eilę ir atsako iš karto;
tada klientas klausia vėl ir vėl, kol darbas baigiasi. Serverio inicijuojamų pranešimų kanalo nėra: pokalbio
srautas yra vienintelis srautas sistemoje. \Cref{fig:api:polling} rodo ataskaitos seką.

\begin{figure}[H]
\centering
\begin{tikzpicture}[node distance=5mm and 6mm]
  \node[svc] (s1) {\code{queued}};
  \node[svc, right=of s1] (s2) {\code{running}};
  \node[man, right=of s2] (s3) {\code{scope_review}};
  \node[svc, right=of s3] (s4) {\code{queued}};
  \node[svc, right=of s4] (s5) {\code{running}};
  \node[svc, right=of s5] (s6) {\code{done}};
  \draw[flow] (s1) -- (s2); \draw[flow] (s2) -- (s3); \draw[flow] (s3) -- (s4);
  \draw[flow] (s4) -- (s5); \draw[flow] (s5) -- (s6);
  \draw[opt] (s2.south) -- ++(0,-5mm) -| (s5.south) node[lbl, pos=.25] {be \code{review}};
  \node[data, above=11mm of s1] (c1) {\code{POST /api/reports}};
  \node[data, above=11mm of s3] (c3) {\code{POST .../scope}};
  \node[data, above=11mm of s6] (c6) {\code{GET .../export}};
  \node[data] (c2) at ($(c1)!.5!(c3)$) {\code{GET} \ldots};
  \node[data] (c5) at ($(c3)!.5!(c6)$) {\code{GET} \ldots};
  \draw[call] (c1) -- (s1) node[lbl, midway] {201};
  \draw[call] (c3) -- (s3) node[lbl, midway] {patvirtinimas};
  \draw[call] (c6) -- (s6) node[lbl, midway] {failas};
  \draw[call] (c2) -- (s2.north);
  \draw[call] (c5) -- (s5.north);
  \node[lanelabel, anchor=south west] at ([yshift=1mm]c1.north west) {kliento užklausos};
  \node[lanelabel, anchor=north west] at ([yshift=-7mm]s1.south west) {ataskaitos \code{status}; ataskaita gali baigtis ir \code{failed} arba \code{cancelled}};
\end{tikzpicture}
\actorlegend
\caption{Ataskaita taip, kaip ją mato klientas. Paleidimas atsako iš karto; klientas periodiškai
klausia \code{GET /api/reports/{id}} apie \code{status} ir \code{progress}; su \code{review}
ataskaita sustoja ir apimties maršrute laukia žmogaus; užbaigtą ataskaitą galima
eksportuoti.}\label{fig:api:polling}
\end{figure}

\paragraph{Ką atsako paleidimas.} Ataskaitos ar programos dokumento paleidimas atsako 201 su
nauja eilute, jau turinčia \code{jobId}, \code{jobStatus} ir \code{progress}. Teisės aktų bazę
keičiantys maršrutai --- atnaujinimas, įkėlimas, pareigų išskyrimas, vertintojas, vektorių
perskaičiavimas, visų šaltinių atnaujinimas --- atsako 202 su pačia užduotimi. Dokumento
įkėlimas yra išimtis kita kryptimi: tekstas išgaunamas prieš atsakymą, ir tik po to vykstantis
indeksavimas yra užduotis, kurios atsakymas neįvardija; ją galima rasti per \route{GET}{/api/jobs}
su \code{subject=document:<id>}.

\paragraph{Dukart ta pati užklausa.} API atmeta arba sujungia antrą paleidimą, užuot tą patį darbą
vykdžiusi du kartus. Antra įmonės ataskaita ar antras tos pačios rūšies dokumentas atsako 409
\code{ALREADY_RUNNING}. Antras to paties teisės akto atnaujinimas, išskyrimas ar vertintojo
paleidimas atsako 202 su užduotimi, kuri jau yra eilėje ar vykdoma, todėl klientas abu atsakymus
gali traktuoti vienodai.

\paragraph{Ką tikrinti.} Ilgo darbo būseną rodo trys ištekliai (\cref{tab:api:polling}).
\code{progress} yra pačios užduoties pranešimas, kur ji yra, --- etapas, procentas ir žinutė.
Užduotis papildomai išvardija savo etapus, kiekvieną su tuo, ką sunaudojo jo modelių
iškvietimai, kai etapas baigiasi.

\begin{table}[H]
\caption{Ištekliai, kuriuos klientas periodiškai tikrina, ir jų būsenos reikšmės}\label{tab:api:polling}
\begin{tabularx}{\textwidth}{@{}L{38mm}Y@{}}
\toprule
\hd{Išteklius} & \hd{\code{status} yra viena iš} \\
\midrule
\code{GET /api/jobs/{id}} & \code{queued}, \code{running}, \code{done}, \code{failed}, \code{cancelled}; su \code{progress}, \code{result}, \code{error}, \code{cancelRequested}, \code{attempts} \\
\code{GET /api/reports/{id}} & \code{queued}, \code{running}, \code{scope_review}, \code{done}, \code{failed}, \code{cancelled}; su \code{jobStatus}, \code{progress}, \code{error} \\
\code{GET /api/programme/{id}} & \code{queued}, \code{running}, \code{draft}, \code{approved}, \code{failed}, \code{cancelled}; su \code{jobStatus}, \code{progress}, \code{error} \\
\bottomrule
\end{tabularx}
\end{table}

Sąrašų maršrutai turi tuos pačius laukus, todėl klientas, rodantis daug eilučių, periodiškai
klausia \route{GET}{/api/reports}, \route{GET}{/api/companies/{id}/programme} ar
\route{GET}{/api/jobs}, o ne kiekvienos eilutės atskirai. Teisės akto išsamioje informacijoje yra
\code{latestJob} --- įkėlimas ar išskyrimas, kuris su juo dirba dabar, arba, jei tokio nėra,
naujausias. Pokalbyje tą patį vaidmenį perkrovus puslapį atlieka asistento paslaugos būsenos
maršrutas (\cref{sec:api:assistant}).

\paragraph{Kur klausiama žmogaus.} Ataskaita, paleista su \code{review: true}, nevyksta iki galo:
nusprendusi, kurie teisės aktai taikomi, ji sustoja būsenoje \code{scope_review}, ir jokia
užduotis su ja nedirba. Ji tęsiama tik tada, kai \route{POST}{/api/reports/{id}/scope} iškviečiamas
su \code{approved: true}, ir tai vėl įtraukia užduotį į eilę. Tą patį maršrutą be
\code{approved} galima kviesti kiek norima kartų pasirinkimui įkainoti; jis nieko neįrašo.
Programos dokumento rengimas baigiasi būsena \code{draft}, o \code{approved} dokumentas tampa tik
žmogui iškvietus \route{PATCH}{/api/programme/{id}}.

\paragraph{Atšaukimas.} Kiekvienas iš trijų išteklių turi atšaukimo maršrutą. Užduotis, ataskaita
ar dokumentas, vis dar esantys eilėje, atšaukiami iš karto; vykdomas gauna žymą, kurios užduočių
vykdytojas paiso per kitą savo patikrą, todėl būsena pasikeičia šiek tiek vėliau; apimties
patvirtinimo laukianti ataskaita atšaukiama iš karto. Viskas, kas jau baigta, atsako 409
\code{NOT_CANCELLABLE}. Ataskaitos ar dokumento, kuris rengiamas, šalinimas atmetamas su 409
\code{STILL_RUNNING}: pirmiausia jį reikia atšaukti.

\paragraph{Pinigai prieš darbą.} Klientas gali paklausti prieš panaudodamas naudotojo biudžetą:
\route{GET}{/api/reports/estimate} pateikia viršutinę ataskaitos kainos ribą kartu su šiandien
likusiu biudžetu, apimties maršrutas įkainoja antrąją dalį pasirinktiems teisės aktams, o
\route{GET}{/api/usage} pateikia šiandienos \code{budget}. Pradėtas vykdymas niekada nestabdomas
pusiaukelėje dėl biudžeto stokos; biudžetas lemia tik tai, ar jį galima pradėti.

\subsection{Failai}\label{sec:api:files}

Trys maršrutai vietoj JSON atsako failu, sukurtu užklausos metu (\cref{tab:api:exports}).
Kiekvienas nustato \code{Content-Disposition: attachment} su failo vardu ir
\code{Cache-Control: no-store}, ir kiekvienas įrašomas į veiksmų žurnalą. Failo varde įmonės pavadinimo
lietuviškos raidės paverčiamos paprastomis lotyniškomis, kiekviena kita seka simbolių, kurie
nėra raidės, skaitmenys, taškas, pabraukimo brūkšnys ar brūkšnelis, virsta vienu brūkšneliu, o
vardas iki plėtinio sutrumpinamas iki 80 simbolių (ataskaitos atveju --- prieš pridedant priešdėlį
\code{legalops-}). Klaida šiuose maršrutuose vis tiek yra JSON apvalkalas, todėl klientas
patikrina būseną prieš traktuodamas turinį kaip failą.

\begin{table}[H]
\caption{Maršrutai, kurie atsako failu}\label{tab:api:exports}
\begin{tabularx}{\textwidth}{@{}L{47mm}L{24mm}Y@{}}
\toprule
\hd{Maršrutas} & \hd{Formatai} & \hd{Failo vardas; sąlyga} \\
\midrule
\code{GET /api/reports/{id}/export} & \code{pdf}, \code{xlsx} & \code{legalops-<company>-<date>}; tik \code{done} ataskaita \\
\code{GET /api/companies/{id}/actions/export} & \code{xlsx} & \code{actions-<company>-<date>}; visada \\
\code{GET /api/programme/{id}/export} & \code{pdf}; \code{docx} (politika); \code{xlsx} (rizikos) & \code{<kind>-<company>-v<version>}; \code{draft} arba \code{approved} dokumentas \\
\bottomrule
\end{tabularx}
\end{table}

Valdymo struktūros dokumento organizacinė schema (\route{GET}{/api/programme/{id}/chart}) yra
vienintelis API teikiamas brėžinys: SVG paveikslas, siunčiamas įterptas, o ne kaip priedas. Failai
į sistemą patenka tik dviem maršrutais, abu \code{multipart/form-data} su dalimi \code{file}:
\route{POST}{/api/documents} --- įmonės dokumentui ir
\route{POST}{/api/corpus/documents/{id}/upload} --- naujai teisės akto redakcijai. Visų kitų
užklausų turinys yra JSON.
